% Modernised reading — fonds Alexandre Grothendieck, shelfmark 85
% (pages 1-97, the whole folder).
%
% This is an INTERPRETATION, not a transcription. It re-reads the pages in
% today's notation and vocabulary, states the hypotheses the manuscript leaves
% implicit, and drops the critical apparatus so that the mathematics reads
% straight through. Everything the brackets and underlines carried moves into a
% footnote. It is held to being mathematically correct as it stands; where the
% manuscript is loose or mistaken, the true statement is what appears here and
% the footnote says what the page has. In any dispute of substance the
% transcription governs: whoever cites Grothendieck must cite batch-01.fr.tex
% (pp. 1-20), batch-02.fr.tex (pp. 21-40), batch-03.fr.tex (pp. 41-60),
% batch-04.fr.tex (pp. 61-80) or batch-05.fr.tex (pp. 81-97).
%
% Pass: Opus 5.5 (claude-opus-5-5), 2026-10-10 — first pass, unchecked against the pages by a human.
% The folder was read whole, its five batches together; this is one document
% for the shelfmark. It was made on Opus 5.5 and has not been measured against
% the reference reading of folder 115 (made on Opus 5). The literature named
% in the footnotes is cited from memory and was not checked against the
% sources. The counts of generating triples (A4: 96, S4: 216, A5: 2280, i.e.
% 38 classes under A5), of subgroups of S5 (156 in 19 classes) and of the
% subgroup classes of D4, S3 x S2, S4, A4, A5, and the roots of the equation
% (E) of pp. 61-73 for each type, were recomputed by machine for this reading.
%
% The transcription is thin. About eighteen \ill/\uncertain per page on
% average; the prose is largely lost on pp. 2, 12-21 (the run on G_1), 44,
% 47, 49, 84-89 and on the scratch sheets 91-93 and 95. Where the prose is
% lost the reading restates the formulas and says so; it never fills a gap.
%
% Scope. Nine covers (pp. 3, 11, 22, 32, 41, 57, 74, 81, 94), each with his
% title and a pencilled page count that matches the leaves: no \pagerange
% claims them, their words head the sections. 88 written pages. Page 96 is
% drawings only. No date anywhere in his hand.
%
% Spine. One group, PGL_2 (his GP(1)), over an arbitrary base, and the regular
% polygons and polyhedra it realises, read through one invariant A (the trace
% 2 cos theta of a rotation):
%   pp. 1-2     work plan: degenerate representations, maps, Schwarz systems,
%               hyperpolygons, Gal(Qbar/Q).
%   pp. 4-10    representations of G_2 in forms of PGL_2: rigid, classified by
%               (alpha_0, alpha_1) with (alpha_0+2)(alpha_1+2)(alpha_0+alpha_1)
%               invertible; Gamma_{nu0,nu1}; S4, S4 x Z/2, A5 x Z/2.
%   pp. 12-21   degenerate (dihedral) case: subgroups of D_inf, faithful D_i,
%               representations of D_inf classified by (alpha, torsor under
%               U_rho), D_i x Z/2.
%   pp. 23-31   the invariant A itself: definition, conjugacy, S_n, orders,
%               F_n, Psi_n, orders in SL_2.
%   pp. 33-40   finite subgroups of PGL_2 in char. 0, normalisers, gerbes,
%               models over Q, forms of D_n, A5 over Spec Z[T]/(T^2+T-1).
%   pp. 42-56   Schwarz systems (generating triples with product 1) of A4, S4,
%               A5; the tetrahedron and the octahedron.
%   pp. 58-73   regular « hyperpolygons » (two pseudo-reflections of C^2): the
%               equation (E) for xi, the census of 44 types, the test.
%   pp. 75-80   finite regular 2-polyhedra in char. 0: nine types.
%   pp. 82-93   Aut of the free group on l_0, l_1, l_2 (l_0 l_1 l_2 = 1), the
%               « eliminations », action on the 96 systems of A4.
%   pp. 95-97   conjugacy classes of subgroups of S5.
%
% Notation fixed for the whole document. PGL_2 for his GP(1); G a form of
% PGL_2 (or of SL_2 where said) over a scheme S; A : G -> A^1_S the invariant
% (his E^1_S for the affine line); alpha = A(u), beta = alpha + 2. D_n the
% dihedral group of order 2n (his double-struck D), D_inf = D_0; S_n, A_n
% symmetric and alternating (his circled S, gothic A). G_1 = D_inf,
% G_2 = <tau_0, tau_1, tau_2 | tau_i^2, (tau_0 tau_2)^2> the cartographic
% groups; rho_0 = tau_0 tau_1 = rho_f, rho_1 = tau_1 tau_2 = rho_s.
% c_5 = 2 cos(2 pi/5), c_5' = 2 cos(4 pi/5), the roots of T^2 + T - 1 (his
% alpha, alpha'; renamed because alpha is the invariant). S_n in Z[U] with
% S_n(T + T^{-1}) = T^n + T^{-n} (his G_n, p. 27; S_n as in folder 78).
% F_n, Psi_n as on p. 29 (= folder 78's). rho with no index is a
% discriminant: 2 - alpha for a polygon (pp. 17-20), -2(alpha_0+alpha_1) for
% a polyhedron (p. 75), as in folder 72. The « eliminations » of pp. 83-93,
% his rho_{eps,eps'.}, rho^I, rho_sigma, are written e here, to keep rho for
% rotations and Schwarz triples. Kappa (p. 61) is A of the product
% tau_0^C tau_1^C. « Schwarz » for his « Schwarz » (p. 1) / « Schwartz »
% (pp. 41 onwards).
%
% Substantive departures from the pages, with page numbers:
%  - p. 6: the product « (alpha_0+2)(alpha_0-2)(alpha_0+alpha_1) » is the
%    (alpha_0+2)(alpha_1+2)(alpha_0+alpha_1) of pp. 4 and 7.
%  - p. 12, D): « two of the three pairs commute » is « one of them » (his
%    i.e.); the equivalences are restated, the proof sketch is mostly lost.
%  - p. 19, Cor. 2 c) (char. 2) describes an empty case: in C, alpha + 2 =
%    4 - rho must be invertible, so rho = 0 forces 2 invertible.
%  - p. 20: U_{rho = alpha - 2} read rho = 2 - alpha (pp. 17, 19).
%  - p. 21, b): D_inf x Z/2 never embeds in PGL_2 in char. 2 (supplied);
%    c): for i >= 4 even, D_i x Z/2 never embeds.
%  - p. 25: « conjugués si G forme de GP(1) » holds over a field, not over a
%    general base (torsors under a torus).
%  - p. 28: the converse (A(u) = -2 => u^2 = 1) holds for forms of PGL_2
%    over reduced S, not for SL_2 (-u unipotent).
%  - p. 29, 2°) and the margin: n invertible on S is needed (Psi_15(-1) = 0
%    mod 5). p. 30, Cor. 1: « resp. 6 » needs residue characteristics != 2.
%  - p. 31, Theorem: A) is « n neither p nor 2p »; B), C) need u nowhere
%    central (supplied, from the margin). With it, the implications <= of
%    A), B), C) hold over any base (Chebyshev argument, given); => holds on
%    reduced bases; on non-reduced bases => is left unchecked.
%  - p. 33: T = Z(Z/n) needs n >= 3. p. 35: N(D_n) = D_2n for every n >= 3
%    (the page: D_n for n even); hence D_2n-torsors in the Cor. Main.
%  - p. 36: « alors G est néc. non tordue » is false (SO(3) over R).
%    p. 37: A5 in PGL_2(k) iff c_5 in k and -1 a sum of two squares, not
%    « k contains the fifth roots of unity ».
%  - p. 38: e) needs n >= 3; 2°) the product on Gamma^- x Gamma^- is the
%    torsor difference, not a constant map.
%  - p. 42 (and pp. 66-67): the table of types omits (3,3,5), (3,3,5'): A5
%    has 38 classes of Schwarz systems under Inn, 19 under Aut, not 32, 16.
%  - p. 48: « 3 · Aut^+ » is Z/2 x Aut^+ (octahedron, icosahedron); p. 49
%    b): involutions, not « d'ordre 3 ».
%  - p. 59: the eigenvalue equations are lambda^2 - x = 0, lambda^2 - y z = 0.
%  - pp. 61-62: the constant term of (E) has 1/16, not 1/8 (his own
%    discriminant requires it). p. 64: lambda^2 - kappa lambda + 1.
%  - pp. 71-72: the trigonometric values confuse pi/5 and 2 pi/5; the
%    conclusions « OK » stand (checked numerically); for (5,3,5'), (3,5,5')
%    the inequality holds, against his « pas vrai ! ».
%  - p. 67: the row (5,5,3), (5,3,5), (3,5,5) counts 3, written « +4 ».
%  - p. 77-78: the image of the tetrahedral G_0 = S4 in PGL_2 is S4, so its
%    signature is (2,3,4), not (2,3,3); p. 78 « A4 » (doubtful) is S4.
%  - p. 80: « OPS alpha_1 > 0 » is not without loss of generality; (4,5'),
%    (5',4), (5',5') pass the definiteness test and are excluded by the
%    finite subgroups of SO(3) (for (5',5'): A5 has no generating pair of
%    5-elements of one class with product of order 2).
%  - pp. 95-97: S5 has 156 subgroups in 19 classes; the page misses the
%    Frobenius group F_20 and counts the pentagons as subgroups (2:1).
%
% Announced and never established in the folder: the « généralités » on
% regular 2-polyhedra invoked on p. 5 (they are folders 70 and 72); the
% inverse functor and torsor structure of p. 18, in the part lost; the
% proposition of p. 20 beyond its statement; the elimination of (5,3,5'),
% (3,5,5') (p. 73); the injectivity questions of pp. 86-87; the finite
% quotients of S3 . E_0 acting on Schwarz systems (p. 89); the outer
% automorphisms of S_n (margin p. 95). Items (2) and (4) of the plan of p. 1
% (Schwarz systems of G_5, Gal(Qbar/Q)) are not reached.
\documentclass[11pt,a4paper]{article}
\input{../preamble/grothendieck}

\folder{85}
\pages{1}{97}
\dating{s.d. — le groupe « Géométrie et topologie combinatoire » (69 à 102) est daté 1976-[vers 1986]}
\watermark{Édition de démonstration\\interprétation personnelle de l'œuvre}
\foldertitle{2-polyèdres et 1-polygônes : notes manuscrites (s.d.) — lecture modernisée du dossier entier}

\begin{document}

\section*{Résumé}

\begin{resume}
Une rotation du plan se reconnaît à un seul nombre : son angle $\theta$, ou,
ce qui revient au même, le nombre $2\cos\theta$, qui est la trace de sa
matrice. Deux rotations de même trace sont conjuguées ; une rotation est
d'ordre $5$ exactement quand $2\cos\theta$ est l'une des deux racines de
$x^2 + x - 1$. Ces pages prennent cette remarque d'école et la poussent aussi
loin qu'elle va. Le groupe qui les occupe est celui des transformations de la
droite projective, $t \mapsto (at + b)/(ct + d)$, que Grothendieck note
$\mathrm{GP}(1)$ et qu'on note aujourd'hui $\mathrm{PGL}_2$ ; c'est aussi,
déguisé, le groupe des rotations de l'espace. À chacun de ses éléments il
attache un invariant $A$, qui joue le rôle de $2\cos\theta$, et il le fait
non seulement sur les nombres réels ou complexes, mais sur une base
quelconque --- les entiers, un corps fini, un anneau où $2$ n'est pas
inversible.

Le problème posé est celui des \emph{polygones et polyèdres réguliers}
réalisés dans ce groupe. Un polyèdre régulier est entièrement décrit par ses
symétries : trois réflexions, dont les produits deux à deux sont la rotation
autour d'une face et la rotation autour d'un sommet. Grothendieck montre que
tout est commandé par les deux invariants de ces deux rotations, $\alpha_0$ et
$\alpha_1$ : à un couple de nombres qui satisfait trois conditions
d'inversibilité correspond un polyèdre régulier, et un seul, à isomorphisme
unique près. Le cube, l'octaèdre, l'icosaèdre ne sont alors que des valeurs
particulières de ce couple ; l'icosaèdre, par exemple, existe sur toute base,
y compris en caractéristique $2$ ou $5$, dès qu'on y a adjoint une racine de
$x^2 + x - 1$. Quand les rotations commutent trop --- le cas « dégénéré » ---
le polyèdre s'écrase en un polygone, et l'étude change de nature.

Le dossier vérifie ensuite ce que dit l'invariant dans les cas finis, en
plusieurs reprises. Il dresse la liste des sous-groupes finis du groupe (ceux
des polygones et des solides de Platon), compte les façons d'engendrer les
groupes du tétraèdre, de l'octaèdre et de l'icosaèdre par trois éléments de
produit $1$ --- les « systèmes de Schwarz » ---, étudie les analogues
complexes des polygones réguliers, où les réflexions sont remplacées par des
transformations d'ordre $3$, $4$ ou $5$, et retrouve enfin, d'un même calcul,
les neuf polyèdres réguliers finis : les cinq solides de Platon et les quatre
polyèdres étoilés de Kepler et Poinsot. Les dernières pages font agir sur
les systèmes de Schwarz un groupe d'opérations formelles ; on y reconnaît
aujourd'hui un premier pas vers l'action du groupe de Galois sur les
« dessins d'enfants » de l'\emph{Esquisse d'un programme}, groupe que le plan
de la première page nomme déjà.

Les feuillets sont des notes de travail, souvent écrites vite et raturées :
plusieurs passages ne se lisent pas, et la lecture qui suit le signale chaque
fois. Grothendieck dit d'ailleurs lui-même, à la page 6, qu'il a « un peu
salopé la présentation », et il laisse des points ouverts qu'il marque
« à vérifier » ou « peut-être faux ». On cherchera ensuite, sous les noms
d'aujourd'hui, la classification de Klein des sous-groupes finis de
$\mathrm{PGL}_2$, les polygones réguliers complexes de Coxeter et les groupes
de réflexions complexes de Shephard et Todd, les polyèdres de Kepler--Poinsot,
et l'action de $\mathrm{Aut}(F_2)$ sur les couples générateurs d'un groupe fini.

\keywords{PGL2 over a base scheme, adjoint quotient, trace invariant, conjugacy classes in PGL2 and SL2, Chebyshev polynomials, minimal polynomial of 2cos(2pi/n), cartographic group, regular map, string Coxeter group, infinite dihedral group, rigid moduli problem, finite subgroups of PGL2, Klein's classification, normalizer of dihedral group, gerbe, generalized dihedral group, icosahedral subgroup of PGL2, generating triples, Schwarz triangle, Nielsen equivalence, regular complex polygon, Shephard-Todd group, complex reflection group, Hermitian form, Kepler-Poinsot polyhedra, Hurwitz formula, automorphisms of free groups, Out(F2), Galois action on dessins, subgroup lattice of S5, Frobenius group F20}
\end{resume}

\section*{Le fil du dossier, et les conventions}

Le dossier compte quatre-vingt-dix-sept pages ; neuf sont des couvertures
(pages 3, 11, 22, 32, 41, 57, 74, 81, 94) qui portent un titre de sa main et,
au crayon, le nombre de feuillets de la liasse qu'elles ouvrent, toujours
exact. Ces titres sont repris ici comme titres de sections. Rien n'est daté ;
le dossier est compté dans le groupe « Géométrie et topologie combinatoire »,
daté 1976--vers 1986. Il appartient à la même famille que les dossiers 70 et
72 (réalisations géométriques des polyèdres réguliers sur une base
quelconque), 76 et 88 (cartes et groupe cartographique), 78 (polygones
réguliers et polynômes semi-cyclotomiques) et 82
($\mathrm{SO}(3) \simeq \mathrm{PGL}_2$) ; on y renvoie là où ils se
rencontrent.

\textbf{Les stations.}
\begin{itemize}
  \item Pages 1--2 : un plan de travail.
  \item Pages 4--10 : représentations de $G_2$ dans les formes de
    $\mathrm{PGL}_2$, cas non dégénéré ; le théorème de classification par
    $(\alpha_0, \alpha_1)$ et les exemples finis.
  \item Pages 12--21 : le cas dégénéré, c'est-à-dire diédral ; représentations
    de $G_1 = \mathbb{D}_\infty$.
  \item Pages 23--31 : l'invariant $A$ lui-même --- le chapitre de fondations
    que les pages 4 à 21 supposent ; il est logiquement premier.
  \item Pages 33--40 : sous-groupes finis de $\mathrm{PGL}_2$ en
    caractéristique $0$, et le groupe $\mathfrak{A}_5$ sur $\mathbf{Z}$.
  \item Pages 42--56 : systèmes de Schwarz de $\mathfrak{A}_4$,
    $\mathfrak{S}_4$, $\mathfrak{A}_5$.
  \item Pages 58--73 : hyperpolygônes réguliers (deux pseudo-réflexions d'un
    plan complexe).
  \item Pages 75--80 : les $2$-polyèdres réguliers finis en caractéristique
    $0$.
  \item Pages 82--93 : opérations sur les systèmes de Schwarz.
  \item Pages 95--97 : classes de sous-groupes de $\mathfrak{S}_5$.
\end{itemize}

\textbf{Conventions, valables pour tout le dossier.} $S$ est un schéma ; $G$
une forme de $\mathrm{PGL}_2$ sur $S$ (de $\mathrm{SL}_2$ quand on le dit) ;
$A : G \to \mathbb{A}^1_S$ l'invariant des pages 23 et suivantes, et, pour une
section $u$ de $G$, $\alpha = A(u)$, $\beta = \alpha + 2$. $\mathbb{D}_n$ est
le groupe diédral d'ordre $2n$, $\mathbb{D}_\infty = \mathbb{D}_0$ le groupe
diédral infini ; $\mathfrak{S}_n$, $\mathfrak{A}_n$ les groupes symétrique et
alterné. Le \emph{groupe cartographique} de dimension $2$ est
\[
  G_2 = \langle \tau_0, \tau_1, \tau_2 \mid \tau_0^2 = \tau_1^2 = \tau_2^2 =
  (\tau_0\tau_2)^2 = 1 \rangle ,
\]
celui de dimension $1$ est $G_1 = \langle \tau_0, \tau_1 \mid \tau_0^2 =
\tau_1^2 = 1\rangle = \mathbb{D}_\infty$ (dossiers 76 et 88) ; on pose
$\rho_0 = \tau_0\tau_1$ (rotation autour d'une face, $\rho_f$) et $\rho_1 =
\tau_1\tau_2$ (rotation autour d'un sommet, $\rho_s$), et
$\Gamma_{\nu_0,\nu_1}$ est le quotient de $G_2$ par $\rho_0^{\nu_0} =
\rho_1^{\nu_1} = 1$, c'est-à-dire le groupe de Coxeter $[\nu_0, \nu_1]$.
On note $c_5 = 2\cos\frac{2\pi}{5} = \frac{-1 + \sqrt5}{2}$ et $c_5' =
2\cos\frac{4\pi}{5} = \frac{-1 - \sqrt5}{2}$ les deux racines de $T^2 + T -
1$\footnote{Il les note $\alpha$ et $\alpha'$ ; on les renomme parce que
$\alpha$ est, dans tout le dossier, la valeur de l'invariant.}. Les
polynômes de $\mathbf{Z}[U]$ sont ceux du dossier 78 : $S_n(T + T^{-1}) =
T^n + T^{-n}$, $F_n$ dont les racines sont les $\zeta + \zeta^{-1}$ pour
$\zeta^n = 1$, $\zeta \neq \pm1$, et $\Psi_n$ le polynôme minimal de
$2\cos\frac{2\pi}{n}$\footnote{Son $S_n$ s'appelle $G_n$ à la page 27 ;
$F_n$ et $\Psi_n$ sont définis à la page 29, et la définition y est
exactement celle du dossier 78. Elle confirme ce que la lecture du
dossier 82 (page 33) supposait de $\Psi_n$.}. Sans indice, $\rho$ désigne un
\emph{discriminant} au sens du dossier 72 : $2 - \alpha$ pour un polygone
(pages 17 à 20), $-2(\alpha_0 + \alpha_1)$ pour un polyèdre (page 75).

\textbf{Ce que le dossier annonce et n'établit pas.} Les « généralités » de
la page 5 sur les $2$-polyèdres réguliers épinglés (elles sont, en substance,
celles des dossiers 70 et 72) ; la structure de torseur de la page 18, dans
la partie perdue ; l'élimination des types $(5, 3, 5')$ et $(3, 5, 5')$ de la
page 73 ; les questions d'injectivité des pages 86--87 et le calcul des
quotients finis de la page 89. Les rubriques (2) et (4) du plan de la page 1
--- les systèmes de Schwarz « de $G_5$ » et le groupe de Galois
$\mathrm{Gal}(\overline{\mathbf{Q}}/\mathbf{Q})$ --- ne sont pas atteintes.

\pagerange{1}{2}
\section*{I. Un plan de travail (pages 1 et 2)}

La page 1 est un plan, en rubriques numérotées de chiffres entourés, dans
l'ordre (1), (1'), (2), (4), (3)\footnote{L'ordre est celui de la page ; la
rubrique (3) est écrite en dernier, tout en bas, sous un trait.} :
\begin{itemize}
  \item[(1)] les représentations dégénérées : de $G_0 = \mathbf{Z}/2$, de
    $G_1 = \mathbb{D}_\infty$ (c'est-à-dire les représentations fidèles de
    $\mathbb{D}_2$, avec les mots « octaèdre » et « quaternions
    généralisés »), les représentations « simplement dégénérées » de $G_2$
    (celles de $\mathbb{D}_\infty \times \mathbf{Z}/2$, et le genre des cartes
    correspondantes), les « doublement dégénérées » (celles de
    $(\mathbf{Z}/2)^3$) ;
  \item[(1')] les points de ramification ; les groupes et cartes associés aux
    représentations projectives de $G_2$ sur un corps, notamment fini ; les
    revêtements de genre $0$ de la sphère privée de trois points ; les cartes
    finies associées aux « hyperpolygônes » réguliers finis ;
  \item[(2)] terminer par les systèmes de Schwarz « pythagoriciens » (« cas de
    $G_5$ ») ;
  \item[(4)] une structure, avec relations, liée à
    $\mathrm{Gal}(\overline{\mathbf{Q}}/\mathbf{Q})$ ;
  \item[(3)] la détermination des hyperpolygônes pythagoriciens.
\end{itemize}
Les rubriques (1), (1') et (3) sont, dans l'ordre inverse, les sections III,
VI--VII et VIII ci-dessous ; (2) et (4) ne sont pas traitées\footnote{Le
texte de la rubrique (4) est en partie illisible ; seul
$\mathrm{Gal}(\overline{\mathbf{Q}}/\mathbf{Q})$ est sûr. « Riemanniens » (deux
fois) est une lecture incertaine. « Pythagoricien » est, dans tout le dossier,
son mot pour ce qui touche aux groupes finis des polyèdres réguliers --- les
« groupes pythagoriciens » $\mathfrak{A}_4$, $\mathfrak{S}_4$,
$\mathfrak{A}_5$ (page 78).}.

La page 2, verso de la première, ne se lit qu'en partie : seule la formule
passe sûrement. Toute application du groupe
\[
  L = \langle \rho_0, \rho_1, \rho_2 \mid \rho_0\rho_1\rho_2 = 1\rangle
\]
--- un groupe libre de rang $2$ --- dans un groupe fini $\Gamma$ définit une
carte finie, et donc tout homomorphisme de $L$ dans un $\mathrm{PGL}_2(k)$
d'image finie donne une telle carte ; les représentations en question ont
ainsi une interprétation géométrique en termes d'hyperpolygônes réguliers, et
les systèmes de Schwarz pythagoriciens en donnent la liste. On lit encore le
nombre « 44 » de types d'hyperpolygônes, qu'on retrouvera à la
page 67\footnote{La prose de la page 2 est presque entièrement illisible ;
on ne garde que ce que confirment la formule, le plan de la page 1 et la
suite du dossier. L'indice de $L$ surcharge un $2$ ; plus bas la page écrit
$L_2$. La note marginale ajoute les hyperpolygônes « dégénérés »,
correspondant aux représentations de $\mathbb{D}_\infty \times
\mathbf{Z}/2$.}. C'est le point de départ de la section IX.

\pagerange{4}{10}
\section*{II. Représentations projectives de $G_2$, cas non dégénéré (pages 4 à 10)}

\pagerange{4}{4}
\subsection*{Le théorème (page 4)}

Soit $\Gamma$ un groupe engendré par trois involutions $\tau_0, \tau_1,
\tau_2$ telles que $\tau_0\tau_2 = \tau_2\tau_0$, c'est-à-dire un quotient de
$G_2$ ; $\rho_0 = \tau_0\tau_1$, $\rho_1 = \tau_1\tau_2$. Soit $S$ un schéma.

\textbf{Théorème} (page 4). Soit $C$ la catégorie dont les objets sont les
couples $(G, \varphi)$, $G$ une forme de $\mathrm{PGL}_2$ sur $S$ et $\varphi
: \Gamma \to G(S)$ un homomorphisme, tels qu'en tout point $s \in S$
\[
  (*) \qquad \rho_{0,s}^2 \neq 1, \quad \rho_{1,s}^2 \neq 1, \quad
  \rho_{0,s}\rho_{1,s} \neq \rho_{1,s}\rho_{0,s},
\]
et dont les flèches sont les isomorphismes. Soit $C'$ la catégorie discrète
des couples $(\alpha_0, \alpha_1)$ de sections de $\mathcal{O}_S$ tels que
\[
  (\alpha_0 + 2)(\alpha_1 + 2)(\alpha_0 + \alpha_1) \quad \text{soit
  inversible.}
\]
Alors :
\begin{enumerate}
  \item pour tout objet $(G, \varphi)$ de $C$, le couple $\alpha_0 =
    A(\varphi(\rho_0))$, $\alpha_1 = A(\varphi(\rho_1))$ est un objet de $C'$ ;
  \item le foncteur $C \to C'$ ainsi défini est pleinement fidèle : $C$ est
    \emph{rigide} (le centralisateur de $\varphi(\Gamma)$ dans $G$ est
    trivial), et deux objets sont isomorphes si et seulement s'ils ont les
    mêmes invariants ;
  \item si $\Gamma = G_2$, c'est une équivalence : tout couple $(\alpha_0,
    \alpha_1)$ de $C'$ provient d'un $(G, \varphi)$, défini à isomorphisme
    unique près.
\end{enumerate}
En d'autres termes, le champ des représentations non dégénérées de $G_2$ dans
les formes de $\mathrm{PGL}_2$ est représentable par l'ouvert de
$\mathbb{A}^2_{\mathbf{Z}}$ complémentaire des trois droites $\alpha_0 = -2$,
$\alpha_1 = -2$, $\alpha_0 + \alpha_1 = 0$\footnote{Cette reformulation est de
l'édition. La page commence un premier énoncé, biffé, puis une condition b)
« l'image de $\Gamma$ dans $G(s)$ n'est pas un groupe diédral », barrée de
grandes croix et remplacée par la non-commutation de $\rho_0$ et $\rho_1$
(variante marginale, biffée elle aussi : la non-commutation de $\rho_0$ et
$\sigma = \rho_0\rho_1 = \tau_0\tau_2$). Le « si $\Gamma = G_2$ » de 2) et
de 3) est ajouté au-dessus de la ligne ; à 3) on lit « si $G = G_2$ »,
lapsus pour $\Gamma$.}. Le produit $(\alpha_0 + 2)(\alpha_1 + 2)(\alpha_0 +
\alpha_1)$ est, au signe et au facteur près, le discriminant
$\beta_0\beta_1\tilde\rho$ de la conique des pages 75 et suivantes : la
condition dit que la conique invariante est lisse (section VIII).

\pagerange{5}{6}
\subsection*{La démonstration, et l'invariant $2\cos$ (pages 5 et 6)}

La démonstration s'appuie sur la théorie générale des $2$-polyèdres réguliers
épinglés\footnote{« On sait, par les généralités » : les généralités ne sont
pas dans ce dossier. Ce sont, en substance, celles du dossier 70 (pages 3 à
25 et 64 à 71) et du dossier 72 (le « dictionnaire » et les coordonnées
universelles). L'édition ne les vérifie pas ici.} : la condition $(*)$ assure
l'existence d'un $2$-polyèdre régulier épinglé, défini à isomorphisme unique
près, dont la représentation $G_2 \to \Gamma \to G(S)$ est $\varphi$. Il reste
à identifier ses paramètres. Dans le contexte général (dimension $n$, dossier
72), le polyèdre est décrit par des paramètres $\beta_i = \alpha_i + 2$
($0 \leq i \leq n - 1$), et ses vecteurs de réflexion $a_i$ ont pour
« longueurs » et « produits scalaires »
\[
  f(a_i) = \beta_0 \cdots \beta_{i-1}, \qquad f(a_{i+1}) = \beta_0 \cdots
  \beta_i, \qquad \varphi(a_i, a_{i+1}) = -\beta_0 \cdots \beta_i ,
\]
d'où, si l'on peut parler de l'angle $\theta_i$ entre $a_i$ et $a_{i+1}$,
\[
  4\cos^2\theta_i = \frac{\varphi(a_i, a_{i+1})^2}{f(a_i)\, f(a_{i+1})} =
  \beta_i, \qquad \text{c'est-à-dire} \qquad \alpha_i = 2(2\cos^2\theta_i - 1)
  = 2\cos 2\theta_i .
\]
Or le produit de deux réflexions d'angle $\theta_i$ est une rotation d'angle
$2\theta_i$, et le lemme suivant dit que l'invariant de la rotation induite sur
la conique est $2\cos 2\theta_i$ : les deux définitions de $\alpha_i$
coïncident\footnote{Une parenthèse « \emph{NB} $\xi_i = \beta_i/(f_i f_{i+1}) =
\beta_i/4 = \cos^2\theta_i$ » suit la formule ; son dénominateur se lit mal et
elle n'est pas utilisée. Dans la marge, un croquis de deux droites, l'angle
$\theta_i$ et une flèche $\rho_i$.}.

\textbf{Lemme} (page 5). Soient $(V, f)$ un module quadratique libre de rang
$3$ sur $S$ et $r$ une rotation d'angle $\vartheta$ autour d'un axe non
isotrope, $\alpha = 2\cos\vartheta$ (c'est-à-dire $\mathrm{Tr}\, r = 1 +
\alpha$). Soit $Y \subset \check{\mathbb{P}}(V)$ la conique des droites
isotropes et $r_Y$ l'automorphisme qu'induit $r$. Alors $A(r_Y) = \alpha$.

\emph{Démonstration.} Localement on peut prendre $f(x, y, z) = x^2 + yz$ et
$r_t(x, y, z) = (x, ty, t^{-1}z)$, de sorte que $\alpha = t +
t^{-1}$\footnote{La page note $\xi$ le paramètre ; on écrit $t$, $\xi$ étant
réservé à la page 61.}. L'homomorphisme $\mathbb{G}_m \to
\mathrm{Aut}(Y)$, $t \mapsto r_t$, est un monomorphisme : la paramétrisation
$(s : s') \mapsto (ss', s^2, -s'^2)$ de $Y$ montre que $r_t$ multiplie le
rapport $s/s'$ par $t$. Par la propriété caractéristique de $A$ (page 23),
$A(r_t) = t + t^{-1} = \alpha$\footnote{Le mot qui qualifie
l'homomorphisme est illisible ; la paramétrisation est de l'édition.}.

Page 6, Grothendieck juge lui-même : « On a un peu salopé la présentation »,
qui omet le cas où les $\beta_i$ ne sont pas inversibles, « sans parler d'une
théorie des angles sur schémas quelconques, que je n'ai pas bien
explicitée ». Mais un principe de prolongement ramène l'énoncé général à la
caractéristique $0$. Comme l'objet de $C$ définit une forme lisse, le produit
$(\alpha_0 + 2)(\alpha_1 + 2)(\alpha_0 + \alpha_1)$ est inversible ;
inversement un couple qui satisfait cette condition définit un polyèdre
régulier, donc une conique lisse et une représentation de
$G_2$\footnote{La page 6 écrit le produit « $(\alpha_0 + 2)(\alpha_0 -
2)(\alpha_0 + \alpha_1)$ », lecture nette ; c'est un lapsus pour celui des
pages 4 et 7. Le premier paragraphe de la page est en grande partie illisible,
et le principe de prolongement est rapporté tel qu'on le devine.}.

\pagerange{6}{8}
\subsection*{Ordres prescrits : le groupe $\Gamma_{\nu_0,\nu_1}$ (pages 6 à 8)}

\textbf{Corollaire} (pages 6 et 7). Soient $\nu_0, \nu_1 \geq 3$ des entiers
et $\Gamma = \Gamma_{\nu_0, \nu_1}$. Soit $C(\nu_0, \nu_1)$ la sous-catégorie
de $C$ formée des $(G, \varphi)$ tels qu'en tout $s \in S$, $\rho_{0,s}$ soit
d'ordre exactement $\nu_0$ et $\rho_{1,s}$ d'ordre exactement $\nu_1$. Soit
$C'(\nu_0, \nu_1)$ la catégorie discrète des $(\alpha_0, \alpha_1)$ tels que
\begin{itemize}
  \item[a)] $(\alpha_0 + 2)(\alpha_1 + 2)(\alpha_0 + \alpha_1)$ soit
    inversible ;
  \item[b)] $\Psi_{\nu_i}(\alpha_i) = 0$ pour $i = 0, 1$ ;
  \item[c)] $\nu_i$ soit inversible sur $S$ si $\nu_i$ n'est pas premier.
\end{itemize}
Alors le foncteur du théorème induit une équivalence $C(\nu_0, \nu_1) \to
C'(\nu_0, \nu_1)$.

La démonstration consiste à traduire « $\rho_i^{\nu_i} = 1$ et $\rho_{i,s}$
d'ordre exactement $\nu_i$ en tout point » en une condition sur l'invariant ;
c'est la traduction de la page 31 (section IV), sous la forme qu'elle prend
pour $\mathrm{PGL}_2$ : un élément d'ordre fini de $\mathrm{PGL}_2(k)$ est
semi-simple d'ordre premier à $p$, ou unipotent d'ordre $p$, de sorte que
$\Psi_n(\alpha) = 0$ suffit si $n$ est premier, et qu'il faut en outre $n$
inversible sinon\footnote{La page dispose la condition en un tableau à trois
cas, « (cas 1°) $\nu_i$ inv. », « (cas 2°) aucune condition », « (cas 3°)
$\nu_i$ inv. », en face d'une rédaction biffée ; on le lit comme renvoyant
aux cas A), C), B) du théorème de la page 31, ce qui donne c). Une
condition c) sur les $\nu_i$ de la forme $2p$, réécrite plusieurs fois, est
barrée. Sur une base non réduite, l'équivalence dépend de cette traduction,
que Grothendieck marque « à vérifier » ; voir la section IV.}. Une remarque
de la page 8, barrée de deux longs traits, observait que pour $\nu_i = 2p$ les
deux catégories sont vides dès que $S$ a un point de caractéristique $p$ :
$\Psi_{2p} \equiv (U + 2)^{(p-1)/2}$ modulo $p$, ce qui force $\alpha_i =
-2$, et $\mathrm{PGL}_2(k)$ n'a pas d'élément d'ordre $2p$ en caractéristique
$p$\footnote{En marge, en bas, un mot biffé lu avec doute « faux ». Le
contenu de la remarque est exact ; c'est sans doute sa conclusion (exiger
« $p$ premier impair ») qui a été abandonnée au profit de c).}.

\pagerange{8}{10}
\subsection*{Les trois cas finis (pages 8 à 10)}

Les groupes $\Gamma_{\nu_0,\nu_1}$ finis avec $\nu_0, \nu_1 \geq 3$ sont ceux
des solides de Platon.

\textbf{A) $\nu_0 = \nu_1 = 3$, $\Gamma \simeq \mathfrak{S}_4$}
(le tétraèdre). Ici $\alpha_0 = \alpha_1 = -1$ et le produit de a) vaut
$1 \cdot 1 \cdot (-2) = -2$. Donc : si $2$ n'est pas inversible sur $S$, la
catégorie est vide --- en caractéristique $2$ toute représentation se
factorise par $\mathfrak{S}_3$ ; si $2$ est inversible (sur
$\mathbf{Z}[\frac12]$), il existe un couple $(G, \varphi)$ et un seul à
isomorphisme unique près, et $\varphi$ est partout fidèle\footnote{Page 8 :
fin de page très raturée, plusieurs $\mathfrak{S}_3$, $\mathbb{D}_3$,
$\mathfrak{S}_4$ biffés ou surchargés. Page 9, l'énoncé b) est serré
d'insertions en partie illisibles ; « partout fidèle » est une lecture
incertaine.}.

\textbf{B) $(\nu_0, \nu_1) = (4, 3)$ ou $(3, 4)$, $\Gamma \simeq \mathfrak{S}_4
\times \mathbf{Z}/2$} (le cube, l'octaèdre). Ici $\{\alpha_0, \alpha_1\} =
\{0, -1\}$ et le produit vaut $-2$ ; il faut encore $2$ inversible. En
caractéristique $2$, $\varphi$ se factorise par $\mathbf{Z}/2 \times
\mathfrak{S}_3 \simeq \mathbb{D}_6$. Sur $S$ où $2$ est inversible, il y a
un couple et un seul, et $\varphi$ est triviale sur le centre $\mathbf{Z}/2$
de $\Gamma$ et induit en chaque point un monomorphisme $\mathfrak{S}_4
\hookrightarrow G_s$\footnote{Le centre de $\Gamma$ est l'antipodie ; elle
agit trivialement sur la conique. Le cas a) de la page 9 est barré et ne
subsiste que par sa note marginale.}.

\textbf{C) $(\nu_0, \nu_1) = (3, 5)$, $\Gamma \simeq \mathbf{Z}/2 \times
\mathfrak{A}_5$} (l'icosaèdre). Ici $\alpha_0 = -1$ et $\alpha_1 = \alpha$
avec $\alpha^2 + \alpha - 1 = 0$. Tout est inversible sans condition :
$\alpha_0 + 2 = 1$, $\alpha + 2$ est de norme $1$ et d'inverse $1 - \alpha$,
$\alpha_0 + \alpha_1 = \alpha - 1$ est de norme $1$ et d'inverse $-\alpha -
2$.

\textbf{Théorème} (page 10). La catégorie des couples $(G, \varphi)$, $G$ forme
de $\mathrm{PGL}_2$ sur $S$ et $\varphi : \mathbf{Z}/2 \times \mathfrak{A}_5
\to G(S)$ non triviale sur $\mathfrak{A}_5$ en chaque point, est équivalente,
par $\varphi \mapsto \alpha(\varphi) = A(\varphi(\rho_1))$, à la catégorie
discrète des $\alpha \in \Gamma(S, \mathcal{O}_S)$ tels que $\alpha^2 + \alpha
- 1 = 0$. Les $\varphi$ sont triviales sur $\mathbf{Z}/2$ et sont donc aussi
les $\varphi^+ : \mathfrak{A}_5 \to G$ non triviales sur toute fibre ---
c'est-à-dire, $\mathfrak{A}_5$ étant simple, injectives sur toute fibre.

La note marginale dit : « laisser tomber $\mathbf{Z}/2$ dans l'énoncé ». C'est
ce que fait le théorème de la page 40 (section V), où ce résultat devient :
le foncteur des $\mathfrak{A}_5$ plongés dans une forme de $\mathrm{PGL}_2$
est représenté par $\mathrm{Spec}\,\mathbf{Z}[T]/(T^2 + T - 1)$. En
particulier l'icosaèdre existe en caractéristique $2$ ($\mathfrak{A}_5 \simeq
\mathrm{PSL}_2(\mathbf{F}_4)$), $3$ ($\mathfrak{A}_5 \subset
\mathrm{PGL}_2(\mathbf{F}_9)$) et $5$ ($\mathfrak{A}_5 \simeq
\mathrm{PSL}_2(\mathbf{F}_5)$, où $T^2 + T - 1 = (T - 2)^2$ et la rotation
d'ordre $5$ est unipotente)\footnote{Ces trois exemples sont de l'édition.}.

\pagerange{12}{21}
\section*{III. Représentations projectives de $G_1$, et représentations dégénérées de $G_2$ (pages 12 à 21)}

Les dix pages de cette liasse sont écrites vite, sous des ratures en
plusieurs couches ; la prose de liaison est souvent perdue et les
démonstrations ne se lisent que par fragments. On donne les énoncés, qui
passent, et ce qu'on lit des preuves.

\pagerange{12}{13}
\subsection*{Les représentations dégénérées (pages 12 et 13)}

\textbf{Proposition} (page 12). Soient $k$ un corps algébriquement clos,
$X$ une droite projective sur $k$, $G = \mathrm{Aut}(X)$, et $\tau_0, \tau_1,
\tau_2 \in G(k)$ des involutions avec $\tau_0\tau_2 = \tau_2\tau_0$. Les
conditions suivantes sont équivalentes :
\begin{itemize}
  \item[A)] il existe un diviseur de degré $2$ de $X$ stable par les $\tau_i$ ;
  \item[A')] les prolongements $\tau_i^C$ au plan projectif $C$ dont $X$ est la
    conique laissent une droite invariante (en caractéristique $\neq 2$ : un
    point invariant) ;
  \item[B)] les $\tau_i$ sont dans un même sous-groupe $H$ de $G$ égal au
    normalisateur $N(T)$ d'un tore maximal ou à un sous-groupe de Borel ;
  \item[C')] $\varphi : G_2 \to G(k)$ se factorise par un quotient de
    $\mathbb{D}_\infty \times \mathbf{Z}/2$ ;
  \item[C'')] $\varphi$ se factorise par une extension d'un groupe commutatif
    par un groupe cyclique ;
  \item[D)] l'un des trois couples $(\tau_0, \tau_1)$, $(\tau_1, \tau_2)$,
    $(\rho_0, \rho_1)$ commute, c'est-à-dire $\rho_0^2 = 1$ ou $\rho_1^2 = 1$
    ou $\rho_0\rho_1 = \rho_1\rho_0$.
\end{itemize}
Ce sont les représentations \emph{dégénérées}, ou \emph{diédrales} : le
groupe engendré est alors un $\mathbb{D}_n$ ($1 \leq n \leq \infty$) ou un
$\mathbb{D}_n \times \mathbf{Z}/2$, et l'on n'a plus affaire à un polyèdre
mais à un polygone\footnote{La page écrit « Deux parmi les trois couples
\ldots{} commutent, i.e. $\rho_0^2 = 1$ ou $\rho_1^2 = 1$ ou $\rho_0\rho_1 =
\rho_1\rho_0$ » ; on suit le « i.e. », qui est la condition juste. Une
condition C), « l'image de $G_2^+$ est commutative », est biffée et remplacée
dans la marge par C') et C'') ; une note marginale encadrée ajoute que dans D),
en caractéristique $\neq 2$, on peut remplacer \ldots{} par $\mathbb{D}_n
\times \mathbf{Z}/2$. La liste des groupes engendrés est lue sous une
rature.}.

Ce qu'on lit des démonstrations : A $\Leftrightarrow$ A' par le
dictionnaire entre $X$ et $C$ (diviseurs de degré $2$ de $X$ et droites de
$C$ --- c'est le dictionnaire du dossier 82) ; A $\Leftrightarrow$ B par les
stabilisateurs des diviseurs de degré $2$, qui sont $N(T)$ (diviseur réduit)
et $B$ (point double) ; B $\Rightarrow$ C, en utilisant que $H(k)$ est un
produit semi-direct $\gamma \cdot U$, $U$ commutatif, et que les $\tau_i$
ont même image dans $\gamma$ quand $(\tau_0, \tau_1)$ et $(\tau_1, \tau_2)$
ne commutent pas, de sorte que $\rho_0, \rho_1 \in U$ commutent ; C
$\Rightarrow$ C' parce que les quotients de $G_2$ par l'une des trois
relations de D) sont tous isomorphes à $\mathbf{Z}/2 \times \mathbb{D}_\infty$ ;
C' $\Rightarrow$ C'' par le sous-groupe cyclique $\mathbb{D}_\infty^+$ ; C''
$\Rightarrow$ A, le diviseur des points fixes d'un générateur $u$ de ce
sous-groupe cyclique étant stable par son normalisateur\footnote{Le
bas de la page 12 et le haut de la page 13 sont lourdement raturés ; les
étiquettes « $D \Rightarrow C''$ » et « B) \& C') » de la dernière phrase
sont transcrites telles qu'elles se présentent et ne s'accordent pas
clairement avec ce qu'elles introduisent. Une marge porte « réf ? ».}.

\pagerange{14}{15}
\subsection*{Les sous-groupes de $\mathbb{D}_\infty$ (pages 14 et 15)}

\textbf{Lemme} (page 14). Soit $\mathbb{D}_\infty = \langle \sigma_0, \sigma_1
\mid \sigma_0^2 = \sigma_1^2 = 1\rangle$, $u = \sigma_0\sigma_1$.
\begin{itemize}
  \item[a)] Pour $i \geq 0$, le sous-groupe $N_i$ engendré par $u^i$ est
    distingué, cyclique infini si $i \neq 0$, trivial si $i = 0$ ; le quotient
    est $\mathbb{D}_i = \langle \sigma_0, \sigma_1 \mid \sigma_0^2 = \sigma_1^2
    = (\sigma_0\sigma_1)^i = 1\rangle$ (avec $\mathbb{D}_0 = \mathbb{D}_\infty$).
  \item[b)] Pour $i \geq 0$ et $j \in \mathbf{Z}/i\mathbf{Z}$, le sous-groupe
    $N_{i,j}$ engendré par $u^i$ et $\sigma_0 u^{j'}$ ($j' \equiv j$) est
    isomorphe à $\mathbb{D}_\infty$ si $i \neq 0$, à $\mathbf{Z}/2$ si $i =
    0$ ; il est distingué si et seulement si $\mathbb{D}_i$ est commutatif,
    c'est-à-dire $i \in \{1, 2\}$ ; il contient un conjugué de $\sigma_0$ si
    et seulement si $j \in 2\mathbf{Z}/i\mathbf{Z}$, un conjugué de $\sigma_1$
    si et seulement si $j + 1 \in 2\mathbf{Z}/i\mathbf{Z}$.
  \item[c)] Les $N_i$, $N_{i,j}$ sont deux à deux distincts, et ce sont tous
    les sous-groupes de $\mathbb{D}_\infty$. Par exemple $\{1\} = N_0$,
    $\mathbb{D}_\infty = N_{1,0}$.
\end{itemize}
\textbf{Corollaire 1.} Les quotients de $\mathbb{D}_\infty$ sont, à
isomorphisme près, les $\mathbb{D}_i$ ($1 \leq i \leq \infty$) et le groupe
trivial. \textbf{Corollaire 2.} Les quotients de $\mathbb{D}_\infty \times
\mathbf{Z}/2$ sont les $\mathbb{D}_i \times \mathbf{Z}/2$ et les
$\mathbb{D}_i$ ; et $\mathbb{D}_i \times \mathbf{Z}/2 \simeq \mathbb{D}_{2i}$
pour $i$ impair\footnote{Le début de l'énoncé est très raturé, et la fin de
c) se lit mal (« sauf \ldots »). Une première description des sous-groupes
de $\mathbb{D}_\infty \times \mathbf{Z}/2$ par un couple $N \subset N'$ de
sous-groupes de $\mathbb{D}_\infty$ est encadrée et barrée. La démonstration du
corollaire 2 (fin de la page 14, haut de la page 15) distingue selon que le
noyau est contenu dans $\mathbb{D}_\infty$ ou non ; elle ne se lit qu'en
partie.}.

\pagerange{15}{16}
\subsection*{Les diédraux de $\mathrm{PGL}_2(k)$ (pages 15 et 16)}

\textbf{Proposition} (page 15). Soient $k$ algébriquement clos de
caractéristique $p$ et $i \in \mathbf{N}$. Il existe un plongement
$\mathbb{D}_i \hookrightarrow \mathrm{PGL}_2(k)$ dans les cas suivants, et
dans ceux-là seulement :
\begin{itemize}
  \item[a)] si $p = 0$, toujours ;
  \item[b)] si $p > 0$ et $i = 0$, si et seulement si $k$ n'est pas algébrique
    sur $\mathbf{F}_p$, c'est-à-dire si $k^\times$ a un élément d'ordre
    infini ;
  \item[c)] si $p > 0$ et $i \geq 2$, si et seulement si $p \nmid i$ ou $i =
    p$.
\end{itemize}
La raison : les éléments d'ordre fini de $\mathrm{PGL}_2(k)$ sont les
semi-simples, d'ordre premier à $p$, et les unipotents $\neq 1$, d'ordre $p$
(page 16) ; dans l'autre sens on place $\mathbb{D}_i$ dans le normalisateur
d'un tore, ou, pour $i = p$, dans un Borel, avec le cas particulier
$\mathbb{D}_2 \subset U$ en caractéristique $2$\footnote{La condition « ou
bien $i = 2p$ et $p$ impair » est biffée, à juste titre. Le $i = 1$ ajouté
en tête est trivial. Une question (« les $\mathbb{D}_i \times
\mathbf{Z}/2$ \ldots{} sont-ils plongeables dans un $N(T)$ ou $B$ ? ») et un
corollaire sur $\mathbb{D}_i \times \mathbf{Z}/2$ (page 16) sont barrés de
grandes croix ; ce corollaire est repris, non barré, à la page 21.}.

\pagerange{16}{20}
\subsection*{Représentations de $\mathbb{D}_\infty$ : le groupe $\mathcal{U}_\rho$ (pages 16 à 20)}

On s'intéresse aux représentations $\varphi : \mathbb{D}_\infty \to G(S)$,
$G$ forme de $\mathrm{PGL}_2$ ; on pose $t_i = \varphi(\tau_i)$, $u =
\varphi(\tau_0\tau_1)$, et l'on se borne d'abord aux $\varphi$ telles qu'en
tout point $u_s^2 \neq 1$, c'est-à-dire que $\varphi_s$ ne se factorise pas
par $\mathbb{D}_2$, ou encore que $t_{0,s}$ et $t_{1,s}$ ne commutent pas. Ce
sont les \emph{polygones réguliers} projectifs du dossier 78, et la condition
dit que $\beta = \alpha + 2$ est inversible\footnote{Le passage des pages 16
et 17 qui introduit les points fixes $a_0, a_1$ et les droites $H_0, H_1$ « comme
d'habitude » et la quadrique $Q$ est en grande partie illisible ; on lit que
$\beta = \alpha + 2$ inversible exprime $a_1 \notin H_0$ et rend la quadrique
lisse.}.

Pour une section $\rho$ de $\mathcal{O}_S$, soit $\mathcal{U}_\rho$ le schéma
en groupes lisse de dimension relative $1$ formé des $\lambda$ tels que $1 +
\lambda\rho$ soit inversible, avec la loi
\[
  [\lambda] \cdot [\lambda'] = [\lambda + \lambda' + \lambda\lambda'\rho] .
\]
L'application $\lambda \mapsto 1 + \lambda\rho$ en est un homomorphisme vers
$\mathbb{G}_m$ ; c'est un isomorphisme si $\rho$ est inversible, et pour
$\rho = 0$ la loi est l'addition. C'est le groupe $\mathcal{U}_\rho$ du dossier
72, qui déforme $\mathbb{G}_m$ en $\mathbb{G}_a$\footnote{La page écrit
$U(\rho)$ et $U_\rho$ ; on suit le dossier 72.}.

\textbf{Théorème} (page 17). Soit $C$ la catégorie des $(G, \varphi)$, $G$
forme de $\mathrm{PGL}_2$ sur $S$, $\varphi : \mathbb{D}_\infty \to G(S)$ avec
$u_s^2 \neq 1$ en tout point. Soit $C'$ la catégorie des couples $(\alpha,
R)$, où $\alpha \in \Gamma(S, \mathcal{O}_S)$ est tel que $\alpha + 2$ soit
inversible et $R$ est un torseur sous $\mathcal{U}_{2-\alpha}$. Pour chaque
$\alpha$, il y a une représentation « privilégiée » $(G_\alpha,
\varphi_\alpha)$ d'invariant $\alpha$, avec $G_\alpha = \mathrm{Aut}(X_\alpha,
Q_\alpha)$, et un isomorphisme $\mathcal{U}_{2-\alpha} \xrightarrow{\sim}
\underline{\mathrm{Aut}}(G_\alpha, \varphi_\alpha)$ ; tordre par $R$ donne un
foncteur $C' \to C$, $(\alpha, R) \mapsto (G_\alpha^R, \varphi_\alpha^R)$, et
c'est une équivalence. Le quasi-inverse associe à $(G, \varphi)$ l'invariant
$\alpha$ de $u$ et le faisceau $R = H_1 \setminus \{c, d\}$, torseur sous
$\mathcal{U}_{2-\alpha} \simeq \underline{\mathrm{Aut}}(X, \varphi)$\footnote{Pages
17 et 18. La lecture des données c) (« représentation privilégiée »,
« $\mathfrak{z}\,\mathrm{Ant}$ », sans doute un centralisateur) et du
quasi-inverse est incertaine ; on ne lit pas ce que sont $c$ et $d$. La page
17 écrit « si $R$ est un torseur sous $G_\alpha$ » : c'est sous
$\mathcal{U}_{2-\alpha}$, comme dans l'énoncé de $C'$.}.

On a $\underline{\mathrm{Aut}}(G, \varphi) \simeq {}_2\mathcal{U}_\rho$, le
sous-groupe des points d'ordre divisant $2$, $\rho = 2 - \alpha$ (page 19) :
le centralisateur de $\varphi(\mathbb{D}_\infty)$ est formé des éléments de
l'image de $\mathcal{U}_\rho$ que $t_0$ inverse et fixe à la fois.

\textbf{Corollaire 1.} Si $\rho = 2 - \alpha$ est inversible, $\mathcal{U}_\rho
\xrightarrow{\sim} \mathbb{G}_m$ par $\lambda \mapsto 1 + \lambda\rho$ ; son
image $T$ dans $G$ est l'unique tore maximal contenant $u$, et
$\underline{\mathrm{Aut}}(G, \varphi) = {}_2T \simeq \mu_2$.

\textbf{Corollaire 2.} Si $\rho = 0$, c'est-à-dire $\alpha = 2$,
$\mathcal{U}_\rho \simeq \mathbb{G}_a$, son image $U$ est un sous-groupe
unipotent maximal contenant $u$, et $\underline{\mathrm{Aut}}(G, \varphi) =
\{1\}$, car $2$ est alors inversible\footnote{La page ajoute c) : « si $2
\cdot 1_S = 0$, $\underline{\mathrm{Aut}}(G, \varphi) = U \simeq \mathbb{G}_a$ ».
Ce cas est vide dans $C$ : $\alpha + 2 = 4 - \rho$ doit y être inversible, et
pour $\rho = 0$ cela force $2$ inversible. L'énoncé c) décrit les
automorphismes de la donnée dégénérée, hors de $C$. Les pages 18 et 19
portaient d'abord trois cas 1) à 3), barrés.}.

\textbf{Corollaire 3.} Pour un groupe $\gamma$ et $\Gamma = \mathbb{D}_\infty
\times \gamma$, la catégorie des $(G, \varphi : \Gamma \to G(S))$ à restriction
non dégénérée est équivalente à celle des triples $(\alpha, R, \varphi_0)$,
$\alpha + 2$ inversible, $R$ torseur sous $\mathcal{U}_\rho$, $\varphi_0 :
\gamma \to ({}_2\mathcal{U}_\rho)(S)$ un homomorphisme\footnote{Page 20, la
page écrit $\mathcal{U}_{\rho = \alpha - 2}$ ; les pages 17 et 19 posent $\rho
= 2 - \alpha$, qu'on suit.}.

Enfin la factorisation par un diédral fini (page 20) : pour $n \geq 3$,
$\varphi$ se factorise par un \emph{monomorphisme} $\mathbb{D}_n
\hookrightarrow G$ (c'est-à-dire $u^n = 1$ et $u_s$ d'ordre exactement $n$ en
tout point) si et seulement si $\Psi_n(\alpha) = 0$ et $n$ est inversible sur
$S$ ou premier\footnote{La page : « $n \cdot 1_S$ est inversible, ou $n = p$
premier impair », puisque $n \geq 3$. « ½cyclotomique » est son abréviation
de semi-cyclotomique, et $\Psi_n$ est le polynôme des pages 7 et 29. La
réserve sur les bases non réduites est celle de la section IV.}. C'est
l'énoncé du corollaire des pages 6 et 7 pour les polygones.

\pagerange{21}{21}
\subsection*{$\mathbb{D}_i \times \mathbf{Z}/2$ dans $\mathrm{PGL}_2(k)$ (page 21)}

\textbf{Corollaire} (page 21). Soient $k$ algébriquement clos et $i \geq 0$
pair. Il existe un homomorphisme injectif $\mathbb{D}_i \times \mathbf{Z}/2
\to \mathrm{PGL}_2(k)$ :
\begin{itemize}
  \item[a)] pour $i = 2$ (le groupe $(\mathbf{Z}/2)^3$), si et seulement si
    $\mathrm{car}\, k = 2$ ;
  \item[b)] pour $i = 0$, si et seulement si $\mathrm{car}\, k = 0$, ou
    $\mathrm{car}\, k = p > 2$ et $k$ non algébrique sur
    $\mathbf{F}_p$\footnote{La page omet « $p > 2$ ». En caractéristique $2$
    l'élément $u$ d'ordre infini est semi-simple, son centralisateur est un tore
    $T \simeq \mathbb{G}_m$, et $T$ n'a pas d'élément d'ordre $2$ : il n'y a
    pas de place pour le facteur $\mathbf{Z}/2$. Le corollaire biffé de la
    page 16 avait, lui, « ssi $p \neq 2$ ».} ;
  \item[c)] pour $i \geq 4$, jamais.
\end{itemize}
Pour c), l'argument de la page : $\mathbb{D}_i$ se plonge si $p \nmid i$, et
alors $\alpha \neq 2$, $\rho \neq 0$, le centralisateur
$\mathrm{Aut}(G, \varphi|\mathbb{D}_i) \simeq \mu_2$ est engendré par $u^{i/2}$,
qui est déjà dans $\mathbb{D}_i$ : le facteur $\mathbf{Z}/2$ ne peut s'envoyer
injectivement\footnote{La page envisage aussi un second cas, « $i = p$ », où
$\alpha = 2$ ; il ne se présente pas, $i$ étant pair et $\geq 4$. La phrase
qui suit « $\mu_2$ » se lit avec peine, et la page s'arrête sur « et en
car. 2 ». La conclusion « jamais » est de l'édition ; elle est ce que
l'argument donne.}.

\pagerange{23}{31}
\section*{IV. Invariants d'automorphismes d'une droite projective (pages 23 à 31)}

Cette liasse est le chapitre de fondations que les sections II et III
supposent.

\pagerange{23}{24}
\subsection*{L'invariant $A$ (pages 23 et 24)}

Soit $G$ une forme de $\mathrm{PGL}_2$ ou de $\mathrm{SL}_2$ sur $S$ ---
un schéma en groupes semi-simple de rang $1$. Le foncteur des fonctions
centrales,
\[
  \mathbf{I}(G) = \underline{\mathrm{Hom}}_{\mathrm{sch}}(G, \mathbb{A}^1_S)^{G}
\]
($G$ agissant par automorphismes intérieurs), est représentable, et il y a
un isomorphisme canonique $A : \mathbf{I}(G) \xrightarrow{\sim}
\mathbb{A}^1_S$ ; autrement dit, l'anneau des fonctions centrales est
$\mathcal{O}_S[A]$, et $A : G \to \mathbb{A}^1_S$ est le \emph{quotient
adjoint}. Il est caractérisé par la propriété : pour tout monomorphisme
$i : \mathbb{G}_m \to G$, et après tout changement de base,
\[
  A(i(\lambda)) = \lambda + \lambda^{-1} .
\]
Pour $\mathrm{SL}_2$, $A$ est la trace ; pour $\mathrm{PGL}_2$, c'est
$(\mathrm{Tr}\, v)^2/\det v - 2$ pour un relèvement $v \in \mathrm{GL}_2$
(page 27)\footnote{Les noms « fonctions centrales », « quotient adjoint » sont
de l'édition ; l'énoncé de représentabilité est le sien, avec, ajouté, « vrai
pour $H$ semi-simple ». Que l'anneau des invariants de $\mathrm{GL}_2$ sur
$\mathbf{Z}$ soit engendré par la trace et le déterminant est un résultat
classique (Donkin pour $\mathrm{GL}_n$) ; la page ne le démontre pas.}.

Valeurs remarquables :
\begin{itemize}
  \item $A(1) = 2$ ;
  \item si $\sigma^2 = 1$ et $\sigma_s \neq 1$ en tout point (une « réflexion »
    de $G$ forme de $\mathrm{PGL}_2$), $A(\sigma) = -2$ : clair si $2$ est
    inversible, car $\sigma$ est alors dans un tore maximal, avec $\lambda =
    -1$ ; en général « par un argument de coefficients universels » ;
  \item si $u$ est \emph{unipotent}, c'est-à-dire s'il existe localement pour
    la topologie fpqc un monomorphisme $i : \mathrm{Aff}(1) = \mathbb{G}_m
    \ltimes \mathbb{G}_a \to G$ et une section $t$ de $\mathbb{G}_a$ avec $u =
    i(t)$, alors $A(u) = 2$. Plus généralement $A(i(\lambda, t)) = \lambda +
    \lambda^{-1}$ : la fonction $f(\lambda, t) = A(i(\lambda, t))$ vaut
    $\lambda + \lambda^{-1}$ pour $t = 0$, est invariante par conjugaison, et
    $(0, t)(\lambda, 0)(0, -t) = (\lambda, (\lambda^{-1} - 1)t)$ ; donc
    $f(\lambda, s) = \lambda + \lambda^{-1}$ dès que $\lambda^{-1} - 1$ est
    inversible, et partout par prolongement.
\end{itemize}
La condition $A(u) = 2$ est donc nécessaire pour que $u$ soit unipotent ; elle
est suffisante sur un corps (tout élément est dans un Borel) ; « j'ignore ce
qui en est dans le cas général »\footnote{Le « 2 » de « $A(u) = 2$ », étroit et
appuyé, pourrait se lire 1 ; le calcul impose 2.}.

\pagerange{25}{26}
\subsection*{L'invariant détecte la conjugaison (pages 25 et 26)}

\textbf{Proposition} (page 25). Soient $u, u'$ deux sections de $G$, et
supposons $(A(u) - 2)(A(u) + 2)$ inversible --- c'est-à-dire $u_s^2$ non
unipotent en tout point. Pour que $u$ et $u'$ soient conjugués localement
pour la topologie fpqc (et même étale), il faut et il suffit que $A(u) =
A(u')$. Si $S$ est le spectre d'un corps et $u, u'$ non centraux, la
conclusion vaut sans restriction sur $A(u)$.

\emph{Démonstration.} L'hypothèse place $u$ et $u'$ chacun dans un unique tore
maximal ; par conjugaison des tores et localisation, c'est le même, et c'est
$\mathbb{G}_m$. Alors $u, u'$ sont des sections $\lambda, \lambda'$ de
$\mathbb{G}_m$ avec $\lambda + \lambda^{-1} = \lambda' + \lambda'^{-1} =
\alpha$ : ce sont deux racines de $t^2 - \alpha t + 1 = 0$, de discriminant
$(\alpha - 2)(\alpha + 2)$ inversible ; l'équation définit un revêtement étale
de degré $2$ de $S$, et localement $\lambda' = \lambda$ ou $\lambda' =
\lambda^{-1}$, conjugués par l'involution de $N(T)$\footnote{La page dit
aussi : « (ou encore, conjugués si $G$ forme de $\mathrm{GP}(1)$) ». C'est
vrai sur un corps : deux matrices non scalaires de même polynôme
caractéristique (à un scalaire près) sont conjuguées dans
$\mathrm{GL}_2(k)$. Sur une base quelconque c'est faux : le transporteur de $u$
à $u'$ est un torseur sous le tore centralisateur, qui peut ne pas être
trivial. Sur une base de Dedekind avec $V = \mathcal{O} \oplus L \simeq M
\oplus M^{-1}L$, les éléments $\mathrm{diag}(1, \lambda)$ relativement aux
deux décompositions ont même invariant et ne sont conjugués que si $M^2$ est
trivial. L'exemple est de l'édition. « Revêtement étale de rang 2 » est une
lecture à vérifier.}.

Aux points où $\alpha = \pm 2$ l'énoncé tombe, et la page 26 en donne la raison
par un calcul infinitésimal. Soit $A$ un anneau artinien local de corps
résiduel $k$, $\lambda \equiv -1$, et $\lambda' = \lambda + \mu$ avec
$\mathfrak{m}\mu = 0$. Alors $\lambda'^{-1} = \lambda^{-1} - \mu$ et
$\lambda' + \lambda'^{-1} = \lambda + \lambda^{-1}$ quel que soit $\mu$, alors
que $\lambda' \in \{\lambda, \lambda^{-1}\}$ n'admet que les deux valeurs
$\mu = 0$ et $\mu = \lambda^{-1} - \lambda$. Dès que $k$ a plus de deux
éléments, on trouve donc $\lambda'$ de même invariant que $\lambda$ sans être
égal à $\lambda^{\pm 1}$\footnote{Le signe devant le dernier $\mu$ de la
deuxième ligne et devant $\mu(\lambda^{-1} - \lambda)$ à la troisième est
surchargé ; avec $\lambda^{-2} \equiv 1$ et $\mathfrak{m}\mu = 0$, les signes
retenus sont ceux qui donnent l'égalité des sommes. La note marginale y voit
« une propriété générale des morphismes ramifiés ». Pour conclure à la
non-conjugaison il faut un peu plus que $\lambda' \neq \lambda^{\pm1}$, le
centralisateur de $u$ n'étant pas $N(T)$ en $\alpha = -2$ ; mais sur $A =
k[\varepsilon]$, $\mathrm{diag}(-1 + a\varepsilon, 1)$ ($a \neq 0$) est de trace
non nulle, et aucun conjugué d'un multiple de $\mathrm{diag}(-1, 1)$ ne l'est :
le contre-exemple tient. Cette vérification est de l'édition.}.

\pagerange{27}{28}
\subsection*{Changer de groupe : $S_n$, $\mathrm{SL}_2 \to \mathrm{PGL}_2$, $\mathrm{GL}_2$ (pages 27 et 28)}

\textbf{Remarque} (page 27). Si $j : \mathbb{G}_m \to G$ a pour noyau $\mu_n$,
il se factorise par $\lambda \mapsto \lambda^n$ suivi d'un monomorphisme, et
\[
  A(j(\lambda)) = \lambda^n + \lambda^{-n} = S_n(\lambda + \lambda^{-1}),
\]
$S_n \in \mathbf{Z}[U]$ unitaire de degré $n$, l'unique polynôme tel que
$S_n(T + T^{-1}) = T^n + T^{-n}$ dans $\mathbf{Z}[T, T^{-1}]^{\sigma} =
\mathbf{Z}[U]$, $\sigma(T) = T^{-1}$ ; c'est le polynôme de Tchebychev
normalisé, $S_n(2\cos\theta) = 2\cos n\theta$\footnote{Il le note $G_n$ ; on
suit le dossier 78. Le nom de Tchebychev est de l'édition.}. Exemples :
\begin{itemize}
  \item pour $p : \widetilde{G} \to G = \widetilde{G}/\mu_2$, $\widetilde{G}$
    forme de $\mathrm{SL}_2$, et toute section $u$ de $\widetilde{G}$ :
    $A(u) = \mathrm{Tr}\, u$ et $A(p(u)) = A(u)^2 - 2$ ;
  \item pour $V$ localement libre de rang $2$, $u \in \mathrm{GL}(V)$ et $u_C$
    l'automorphisme induit de $C = \mathbb{P}(V)$ :
    \[
      A(u_C) = \frac{(\mathrm{Tr}\, u)^2}{\det u} - 2 ,
    \]
    ce qui se ramène au cas précédent en écrivant localement $u = tv$, $t$ une
    unité, $v \in \mathrm{SL}(V)$. Si les valeurs propres de $u$ sont
    $\lambda_1, \lambda_2$, $A(u_C) = \lambda_1/\lambda_2 +
    \lambda_2/\lambda_1$ ; pour une \emph{pseudo-réflexion} de déterminant
    $\lambda$ (valeurs propres $\lambda$ et $1$), $A = \lambda +
    \lambda^{-1}$, ce qu'on voit aussi en l'insérant dans la famille à un
    paramètre $\mathbb{G}_m \hookrightarrow \mathrm{GL}(V) \to
    \underline{\mathrm{Aut}}_C$.
\end{itemize}
Ce dernier calcul est celui dont se servira la section VII, où $\kappa$ est
l'invariant du produit de deux pseudo-réflexions.

\pagerange{28}{29}
\subsection*{Les involutions (pages 28 et 29)}

\textbf{Proposition} (page 28). Soit $u$ une section de $G$ avec $u_s \neq 1$
en tout point et $u^2 = 1$ ; si $G$ est une forme de $\mathrm{SL}_2$, on
suppose les caractéristiques résiduelles $\neq 2$. Alors $A(u) = -2$.

Pour une forme de $\mathrm{PGL}_2$, $u$ est dans un unique tore maximal, où il
correspond à $-1$, et le résultat vaut sans restriction de caractéristique
parce que le schéma des réflexions de $\mathrm{PGL}_{2,\mathbf{Z}}$ est lisse
sur $\mathbf{Z}$. Pour $\mathrm{SL}(V)$, Cayley--Hamilton donne $u^2 -
(\mathrm{Tr}\, u) u + 1 = 0$, donc $u^2 = 1$ équivaut à $(\mathrm{Tr}\, u)u =
2$ ; en caractéristique $\neq 2$ cela force $u_s$ scalaire, $u_s = -1$, et
$\mathrm{Tr}\, u = -2$\footnote{La fin de l'argument, en caractéristique
résiduelle $2$ (« Faisant un calcul infinitésimal \ldots »), est biffée et se
perd. La page ajoute entre crochets : « la réciproque étant vraie si $S$ est
réduit ». Pour une forme de $\mathrm{PGL}_2$ sur $S$ réduit, c'est exact.
Pour $\mathrm{SL}_2$, c'est faux : $-v$, $v$ unipotent $\neq 1$, est de trace
$-2$ et d'ordre $2p$ ou infini.}.

\pagerange{29}{31}
\subsection*{L'ordre d'un élément, lu sur son invariant (pages 29 à 31)}

Soient $F_n$ et $\Psi_n \in \mathbf{Z}[U]$ les polynômes unitaires dont les
racines sont les $\zeta + \zeta^{-1}$, $\zeta$ parcourant respectivement les
racines $n$-ièmes de l'unité autres que $\pm 1$ (chaque paire
$\{\zeta, \zeta^{-1}\}$ une fois) et les racines primitives $n$-ièmes ; on a
$F_n = \prod_{d \mid n,\, d \geq 3} \Psi_d$, et $\Psi_n$ est le polynôme minimal
de $2\cos\frac{2\pi}{n}$ (dossier 78).

\textbf{Proposition} (page 29). Soient $G$ une forme de $\mathrm{PGL}_2$ ou de
$\mathrm{SL}_2$, $u$ une section de $G$, $\alpha = A(u)$, $n \geq 3$, et
supposons $(\alpha + 2)(\alpha - 2)$ inversible. Alors :
\begin{itemize}
  \item[1°)] $u^n = 1$ si et seulement si $F_n(\alpha) = 0$ ;
  \item[2°)] $u^n = 1$ et $u_s$ d'ordre exactement $n$ en tout point si et
    seulement si $\Psi_n(\alpha) = 0$ \emph{et $n$ est inversible sur $S$}.
\end{itemize}
Le principe : $u$ est dans un unique tore maximal $T \simeq \mathbb{G}_m$,
correspond à une section $\lambda$, $\alpha = \lambda + \lambda^{-1}$, et
l'on est ramené au calcul : $F_n(\lambda + \lambda^{-1})$ est, au
produit près par une puissance de $\lambda$, $(\lambda^n - 1)/(\lambda^{\epsilon} -
1)$, $\epsilon \in \{1, 2\}$ selon la parité de $n$, où $\lambda^{\epsilon} - 1$ est inversible par
hypothèse\footnote{La condition « $n$ inversible » de 2°) est ajoutée : sans
elle l'énoncé est faux. Modulo $5$, $\Psi_{15} \equiv \Psi_3^4$, de sorte
qu'en $\alpha = -1$, sur $\mathbf{F}_5$, $\Psi_{15}(\alpha) = 0$ et
$\alpha^2 - 4 = -3$ est inversible, alors que $u$ est d'ordre $3$. La note
marginale de la page donne la version sans hypothèse sur $\alpha$ :
« [$u^n = 1$ et $u_s$ d'ordre $n$ premier aux caractéristiques résiduelles]
$\Leftrightarrow$ [$\Psi_n(\alpha) = 0$ et $\alpha^2 - 4$ inversible] » ; il faut
y ajouter, à droite, « $n$ inversible », pour la même raison. Le dernier signe
de 1°) ressemble à un « 1 » ; on lit $F_n(\alpha) = 0$.}.

\textbf{Corollaires} (page 30). Si les caractéristiques résiduelles sont $\neq 3$
(resp. $\neq 2, 3$), $u$ est d'ordre exactement $3$ (resp. $6$) en tout point
si et seulement si $\alpha = -1$ (resp. $\alpha = 1$). Si elles sont $\neq 2$,
d'ordre $4$ si et seulement si $\alpha = 0$. Si elles sont $\neq 5$, d'ordre
$5$ si et seulement si $\alpha^2 + \alpha - 1 = 0$\footnote{La page ne suppose,
pour l'ordre $6$, que les caractéristiques $\neq 3$ ; il faut aussi $\neq 2$ :
en caractéristique $2$, $\alpha = 1 = -1$ donne un élément d'ordre $3$. Pour
l'ordre $5$, la page écrit « $u^5 = 0$ », lapsus pour $u^5 = 1$.}.

\textbf{Les ordres dans $\mathrm{SL}_2(k)$} (pages 30 et 31). L'ordre d'un
élément de $\mathrm{SL}_2(k)$, $k$ un corps de caractéristique $p$, est infini,
ou un entier premier à $p$ (semi-simple), ou $p$ (unipotent $\neq 1$), ou
$2p$ avec $p$ impair ($-1$ fois un unipotent $\neq 1$). Pour $u$ non central,
la valeur de $\alpha$ détermine la classe de conjugaison localement, et, en
notant $\zeta$ une racine de $\zeta^2 - \alpha\zeta + 1$ dans $\bar k$ : en
caractéristique $0$, $u$ est d'ordre fini si et seulement si $\Psi_n(\alpha) =
0$ pour un $n \geq 3$ ; en caractéristique $p > 0$, si et seulement si $\alpha$
est algébrique sur $\mathbf{F}_p$, si et seulement si $\Psi_n(\alpha) = 0$ pour
un $n$ premier à $p$, égal à $p$ ou à $2p$ --- et cet $n$, pris $\geq 3$ ou
égal à $p$, $2p$, est alors l'ordre de $u$\footnote{Il écrit « $n$ est alors
unique, et égal à $\nu(u)$ ». L'unicité demande d'écarter $n = 1, 2$, dont les
racines $\pm 2$ sont aussi celles de $\Psi_p$ et $\Psi_{2p}$ modulo $p$ ; c'est
ce que fait l'hypothèse « $u$ non central », soulignée deux fois dans la
marge. Le second membre de « $\Psi_n(\alpha) = $ » en caractéristique $0$ est
laissé en blanc.}.

\textbf{Théorème} (page 31). Soient $\widetilde{G}$ une forme de
$\mathrm{SL}_2$ sur $S$, $u$ une section de $\widetilde{G}$ \emph{nulle part
centrale}, $\alpha = A(u)$, $n \geq 3$, et
\[
  (*) \qquad u^n = 1 \quad \text{et} \quad u_s \text{ d'ordre exactement } n
  \text{ pour tout } s \in S .
\]
\begin{itemize}
  \item[A)] Si $n$ n'est ni un premier impair $p$, ni le double $2p$ d'un tel
    premier, $(*)$ équivaut à : $\Psi_n(\alpha) = 0$ et $n$ inversible sur $S$.
  \item[B)] Si $n = 2p$, $p$ premier impair, $(*)$ équivaut à :
    $\Psi_n(\alpha) = 0$ et $2$ inversible sur $S$.
  \item[C)] Si $n = p$ premier impair, $(*)$ équivaut à : $\Psi_p(\alpha) = 0$.
\end{itemize}
Ces conditions entraînent $\alpha^2 - 4$ inversible, sauf en B) et C) aux
points de caractéristique $p$, où $\alpha_s = -2$ (cas B) ou $\alpha_s = +2$
(cas C).

Grothendieck marque B) et C) « à vérifier » et C) « peut-être faux ».
L'hypothèse « nulle part central », qui est celle de la marge de la page 31,
les rend justes, en tout cas sur une base réduite. Le sens « $\Leftarrow$ » vaut
même sur toute base, par un argument universel : dans $\mathrm{SL}_2$, pour
tout $n$, $u^n = S'_{n-1}(\alpha)\, u - S'_{n-2}(\alpha)$, où $S'_m$ est le
polynôme de Tchebychev de deuxième espèce normalisé ($S'_m(\lambda +
\lambda^{-1}) = (\lambda^{m+1} - \lambda^{-m-1})/(\lambda - \lambda^{-1})$) ;
$F_n$ divise $S'_{n-1}$ et $S'_{n-2} + 1$ dans $\mathbf{Z}[U]$ (c'est vrai sur
les racines, et $F_n$ est unitaire séparable), donc $F_n(\alpha) = 0$ entraîne
$u^n = 1$ sur toute base. Pour C), $F_p = \Psi_p$ et $\Psi_p(2) = p$ : aux points
de caractéristique $\neq p$, $u_s \neq 1$, et aux points de caractéristique $p$,
$u_s$ est unipotent et non trivial ; il est d'ordre $p$ partout. Pour B), on
remplace $u$ par $-u$, puisque $\Psi_{2p}(U) = \Psi_p(-U)$, et $u_s \neq \pm 1$,
$2$ inversible, donnent l'ordre $2p$ exactement. Pour A), $n$ inversible rend
$\Psi_n$ séparable sur chaque fibre. Le sens « $\Rightarrow$ » se vérifie
fibre à fibre, ce qui suffit sur une base réduite ; sur une base non réduite,
cette lecture ne l'a pas vérifié\footnote{La page énonce A) « Supposons
que $n$ ne soit [pas] de la forme $p$ ($p$ premier impair) », sur une
rédaction biffée « ni de \ldots{} premier, ni $2p$ » ; c'est la seconde qui
convient, B) traitant $2p$. Une note marginale sur $n = 4$ demande qu'A)
« marche dans le cas de $G = \widetilde{G}/\mu_2$ ». L'argument par les
polynômes de Tchebychev est de l'édition. Sans « nulle part central », B) et
C) sont faux : $u = -1$ (resp. $u = 1$) en un point de caractéristique $p$
satisfait les conditions de droite sans être d'ordre $2p$ (resp. $p$).}.

\pagerange{33}{40}
\section*{V. Sous-groupes finis de $\mathrm{PGL}_2$ (pages 33 à 40)}

\pagerange{33}{33}
\subsection*{La liste, en caractéristique $0$ (page 33)}

Soit $k$ algébriquement clos de caractéristique $0$, $G = \mathrm{PGL}_2(k)$.

\textbf{Théorème} (page 33). a) Deux sous-groupes finis de $G$ sont conjugués
si et seulement s'ils sont isomorphes. b) Les types d'isomorphie de
sous-groupes finis de $G$ sont $\mathbf{Z}/n$ ($n \geq 1$), $\mathbb{D}_n$
($n \geq 2$), $\mathfrak{A}_4$, $\mathfrak{S}_4$, $\mathfrak{A}_5$ ---
« exactement les types des groupes d'automorphismes des cartes orientées
régulières sphériques ! ». c) Chacun est son propre normalisateur, sauf
$\mathbf{Z}/n$, $\mathbb{D}_2$, $\mathfrak{A}_4$ (et, voir plus bas,
$\mathbb{D}_n$), avec, pour ceux-là :
\[
  \mathcal{N}(\mathbf{Z}/n) = \mathcal{N}(T), \qquad
  \mathcal{N}(\mathbb{D}_2) \simeq \mathfrak{S}_4, \qquad
  \mathcal{N}(\mathfrak{A}_4) = \mathfrak{S}_4 ,
\]
$\mathcal{N}(\mathbb{D}_2)/\mathfrak{Z}(\mathbb{D}_2) \simeq
\mathrm{Aut}(\mathbb{D}_2) \simeq \mathfrak{S}_3$ et $\mathcal{N}(\mathfrak{A}_4)
= \mathrm{Aut}(\mathfrak{A}_4)$\footnote{La classification est celle de Klein
(\emph{Vorlesungen über das Ikosaeder}, 1884) ; la page ne cite rien. En
caractéristique $p$ elle est remplacée par celle de Dickson. La moitié
inférieure de la page est couverte de traits obliques au crayon, parallèles
et croisés, dont on ne sait s'ils biffent c) ; le texte à l'encre reste
lisible. La parenthèse de b) reste ouverte, et la note sur les cartes
continue dans la marge par des mots non lus. $\mathcal{N}$ et $\mathfrak{Z}$
rendent normalisateur et centralisateur. La réserve sur $\mathbb{D}_n$ est
de l'édition, voir la page 35.}.

Puis, pour toute carte orientée finie sphérique connexe $C$, la réalisation
topologique donne un foncteur de groupoïdes de $C$ vers la catégorie des
« cartes holomorphes sphériques », puis vers celle des couples $(G, \Gamma)$,
$G$ forme de $\mathrm{PGL}_2(\mathbf{C})$, $\Gamma$ sous-groupe fini : une carte
régulière de la sphère, une fois rendue holomorphe, voit son groupe
d'automorphismes devenir un sous-groupe fini de $\mathrm{PGL}_2(\mathbf{C})$.
C'est le lien entre ce dossier et les dossiers 76 et 88.

\pagerange{34}{37}
\subsection*{Normalisateurs, et les catégories de couples $(G, \Gamma)$ (pages 34 à 37)}

Dans chaque cas on décrit le centralisateur $\mathfrak{Z}(\Gamma)$, le
normalisateur $\mathcal{N}(\Gamma)$, qui est le groupe
$\mathrm{Aut}(G, \Gamma)$ des automorphismes du couple (tout automorphisme
d'une forme de $\mathrm{PGL}_2$ étant intérieur), et le quotient
$\mathcal{N}/\mathfrak{Z}$, qui s'envoie dans $\mathrm{Aut}(\Gamma)$.

1) $\Gamma \simeq \mathbf{Z}/n$, $n \geq 2$. $\Gamma$ est contenu dans un
unique tore maximal $T$, et $\Gamma = {}_nT$ ; les applications $\Gamma
\mapsto T$, $T \mapsto {}_nT$ sont des bijections inverses entre tores
maximaux et sous-groupes cycliques d'ordre $n$. Pour $n \geq 3$, $T =
\mathfrak{Z}(\Gamma)$, $\mathcal{N}(\Gamma) = \mathcal{N}(T)$ et
$\mathcal{N}/\mathfrak{Z} \simeq \mathbf{Z}/2$\footnote{La page pose $T =
\mathfrak{Z}(\Gamma)$ dès $n \geq 2$. Pour $n = 2$ le centralisateur d'une
involution est $\mathcal{N}(T)$ tout entier, et $\mathcal{N}/\mathfrak{Z}$
est trivial ; le reste vaut. En marge, deux dessins : une étoile de huit
arêtes issues d'un sommet, un cercle marqué de deux points.}.

2) $\Gamma = \mathbb{D}_n$, $n \geq 3$. Soit $\Gamma^+ \simeq \mathbf{Z}/n$
le sous-groupe d'indice $2$, $T$ l'unique tore maximal qui le contient. Alors
$\mathfrak{Z}(\Gamma) = {}_2T \simeq \{\pm 1\}$, et
\[
  \mathcal{N}(\Gamma) = {}_{2n}T \rtimes \langle\sigma\rangle \simeq
  \mathbb{D}_{2n} ,
\]
quel que soit $n$ : un élément $t$ de $T$ normalise $\Gamma$ si et seulement si
$t\sigma t^{-1} = t^2\sigma \in \Gamma$, c'est-à-dire $t^2 \in {}_nT$. Le
quotient $\mathcal{N}/\mathfrak{Z} \simeq \mathbb{D}_n$ s'envoie dans
$\mathrm{Aut}(\mathbb{D}_n) \simeq \mathrm{Aff}(1, \mathbf{Z}/n)$\footnote{La
page donne $\mathcal{N}(\Gamma) = \Gamma$ si $n$ est pair et $\Gamma \cdot
{}_2T$ si $n$ est impair. Pour $n$ impair c'est bien $\mathbb{D}_{2n}$ ; pour
$n$ pair c'est faux : dans $\mathrm{SO}(3)$, la rotation d'angle $\pi/4$ autour
de l'axe d'un prisme carré normalise son groupe $\mathbb{D}_4$ sans y
appartenir. La dernière formule de la page, « $\mathcal{N}/\mathfrak{Z} \simeq
\Gamma/{}_2\Gamma^+ \subset \mathrm{Aut}(\Gamma)$ », est à corriger de même.
En marge, un « vérifier » (?) encerclé.}.

3) $\Gamma = \mathbb{D}_2 \simeq (\mathbf{Z}/2)^2$. $\mathfrak{Z}(\Gamma) =
\Gamma$, $\mathcal{N}(\Gamma) \simeq \mathfrak{S}_4$ est l'unique sous-groupe
de type $\mathfrak{S}_4$ contenant $\Gamma$, et $\Gamma$ est l'unique
sous-groupe distingué d'ordre $4$ de $\mathcal{N}(\Gamma)$ : d'où une
bijection entre sous-groupes de type $\mathbb{D}_2$ et sous-groupes de type
$\mathfrak{S}_4$ ; $\mathcal{N}/\mathfrak{Z} \simeq \mathrm{Aut}(\Gamma)
\simeq \mathfrak{S}_3$.

4) et 5) $\Gamma \simeq \mathfrak{A}_4$ ou $\mathfrak{S}_4$.
$\mathfrak{Z}(\Gamma) = 1$, $\mathcal{N}(\Gamma) \simeq \mathfrak{S}_4
\xrightarrow{\sim} \mathrm{Aut}(\Gamma)$.

6) $\Gamma \simeq \mathfrak{A}_5$ (page 37). $\mathfrak{Z}(\Gamma) = 1$,
$\mathcal{N}(\Gamma) = \Gamma$, d'indice $2$ dans $\mathrm{Aut}(\Gamma) \simeq
\mathfrak{S}_5$\footnote{Ce cas est d'abord écrit « 5) » à la page 36, encadré
et barré, puis repris page 37.}.

Ces calculs donnent, sur une base $S$ de caractéristique $0$, des
« corollaires principaux » : la catégorie des couples $(G, \Gamma)$, $G$ forme
de $\mathrm{PGL}_2$, $\Gamma \subset G$ forme d'un groupe fini donné, est une
gerbe, la gerbe des torseurs sous le normalisateur du couple standard.
\begin{itemize}
  \item Pour $\Gamma$ forme de $\mathbb{D}_2$, de $\mathfrak{A}_4$ ou de
    $\mathfrak{S}_4$ (pages 34 et 35) : elle équivaut à la catégorie des
    torseurs sous $\mathfrak{S}_4$, ou des revêtements étales de degré $4$ de
    $S$, ou des tétraèdres sur $S$, ou des octaèdres orientés, ou des formes
    de $\mathfrak{A}_4$, ou de $\mathfrak{S}_4$ ; et les sous-groupes de ces
    trois types contenus dans un $G$ fixé se correspondent.
  \item Pour $\Gamma$ forme de $\mathbb{D}_n$, $n \geq 3$ : avec le
    normalisateur corrigé, ce sont les torseurs sous $\mathbb{D}_{2n}$, pour
    tout $n$\footnote{Page 35 : pour $n$ pair, la page conclut aux
    revêtements principaux de groupe $\mathbb{D}_n$, « i.e. celle des
    polygones réguliers sur $S$ de degré $n$ », et pour $n$ impair renvoie à
    $\Gamma' = \mathbb{D}_{2n}$. Avec $\mathcal{N}(\mathbb{D}_n) \simeq
    \mathbb{D}_{2n}$, c'est $\mathbb{D}_{2n}$ dans les deux cas. Le « $\geq 3$ »
    est surchargé, peut-être en « $\geq 4$ ».}.
  \item Pour $\Gamma$ cyclique d'ordre $n \geq 2$, la page conclut à « une
    [\ldots] de rang $2$ sur $S$ », que l'on ne lit pas : le normalisateur est
    $\mathcal{N}(T)$, non fini, et la gerbe est celle des torseurs sous
    $\mathcal{N}(T)$.
\end{itemize}

\textbf{Modèles sur $\mathbf{Q}$} (pages 36 et 37). Chaque gerbe a un
représentant canonique défini sur $\mathbf{Q}$ : $\mu_n$ et $\mathbf{Z}/2
\ltimes \mu_n$ agissant sur $\widehat{\mathbb{G}}_{m,\mathbf{Q}} \simeq
\mathbb{P}^1_{\mathbf{Q}}$ --- tous plongés dans $\mathbf{Z}/2 \ltimes
\mathbb{G}_m$, normalisateur du tore standard de $\mathrm{PGL}_{2,\mathbf{Q}}$ ;
puis le normalisateur dans $\mathrm{PGL}_2$ du $\mathbb{D}_2$ standard, qui
est une forme de $\mathfrak{S}_4$, « (tordue ?) ». Elle l'est : aucun
$\mathrm{PGL}_2(K)$ ne contient $\mathfrak{S}_4$ si $-1$ n'est pas somme de deux
carrés dans $K$\footnote{Critère de Beauville (« Finite subgroups of
$\mathrm{PGL}_2(K)$ », 2010), cité de mémoire : en caractéristique $0$,
$\mathrm{PGL}_2(K)$ contient $\mathbf{Z}/n$ ou $\mathbb{D}_n$ si et seulement si
$2\cos\frac{2\pi}{n} \in K$, contient $\mathfrak{A}_4$ ou $\mathfrak{S}_4$ si et
seulement si $-1$ est somme de deux carrés dans $K$, et $\mathfrak{A}_5$ si et
seulement si de plus $\sqrt5 \in K$. Pour $\mathfrak{S}_4$ constant, on peut
prendre la forme $\mathrm{SO}(x^2 + y^2 + z^2)$ et le groupe des matrices de
permutation signées de déterminant $1$. La note marginale cherche « un autre
représentant plus commode ».}.

Une remarque, reprise en « ReNB » : sur un corps $k$ et pour $n$ donné, sont
équivalents a) $\Psi_n$ a une racine dans $k$ ; b) il existe une forme non
tordue de $\mathbf{Z}/n$ dans une forme de $\mathrm{PGL}_2$ ; c) de même avec
$\mathbb{D}_n$. Sur $\mathbf{Q}$, cela n'arrive que pour $n \in \{2, 3, 4,
6\}$\footnote{La page ajoute en b) : « (alors $G$ est néc. non tordue) ». C'est
faux : $\mathrm{SO}(3)$ sur $\mathbf{R}$, forme non déployée de
$\mathrm{PGL}_2$, contient les rotations de tous les ordres.}.

Pour $\Gamma \simeq \mathfrak{A}_5$ (page 37) : une fois choisi un couple de
référence $(G_0, \Gamma_0)$ sur $\mathbf{Q}$, se donner $(G \supset \Gamma)$
sur $S$ de caractéristique $0$ revient à se donner un torseur sous
$\Gamma_0$, une forme intérieure ; si $\Gamma_0 \simeq \mathfrak{A}_5$
constant, cela revient à un revêtement de degré $5$ de $S$ muni d'une
trivialisation du revêtement quadratique de ses orientations. Il existe sur
$k$ une forme non tordue de $\mathfrak{A}_5$ dans une forme de $\mathrm{PGL}_2$
si et seulement si $T^2 + T - 1$ a ses racines dans $k$\footnote{La page
ajoute, marqué « vérifier » souligné deux fois : « Pour que de plus $G$ soit
non tordue, il f. et s. que $k$ contienne les racines cinquièmes de 1 ».
D'après le critère cité plus haut, la condition juste est que $-1$ soit somme
de deux carrés dans $k$ ; $\mathbf{Q}(\sqrt5, i)$ la satisfait sans contenir
$\mu_5$.}. Sur un $k$ muni d'une racine $\alpha$, un plongement de
$\mathfrak{A}_5$ est défini à isomorphisme unique près par la condition que
l'invariant d'un $5$-cycle standard soit $\alpha$ ; on peut prendre $k = k_0 =
\mathbf{Q}[T]/(T^2 + T - 1)$, et sur les $k_0$-schémas la donnée de $(G \supset
\Gamma)$ équivaut à celle d'un torseur sous $\mathfrak{A}_5$.

\pagerange{38}{39}
\subsection*{Formes de $\mathbb{D}_n$ (pages 38 et 39)}

Soit $\Gamma$ un groupe admettant un sous-groupe $\Gamma^+$ tel que : a)
$\Gamma^+$ est d'indice $2$ ; b) $\Gamma^+$ est commutatif ; c)
$\Gamma/\Gamma^+$ agit sur $\Gamma^+$ par $x \mapsto -x$ ; d) l'extension est
scindée ; e) $\Gamma^+ \setminus {}_2\Gamma^+$ engendre $\Gamma^+$. C'est ce
qu'on appelle aujourd'hui un groupe \emph{diédral généralisé}
$\mathrm{Dih}(\Gamma^+)$. Alors $\Gamma^- = \Gamma \setminus \Gamma^+$ est
formé d'involutions, ${}_2\Gamma = \Gamma^- \amalg {}_2\Gamma^+$, et
\[
  \Gamma^+ = \text{le sous-groupe engendré par } \Gamma \setminus {}_2\Gamma ,
\]
de sorte que $\Gamma^+$ est déterminé par $\Gamma$. De plus $\Gamma^-$ est un
torseur sous $\Gamma^+$, et la multiplication de $\Gamma$ se reconstruit à
partir de celle de $\Gamma^+$ et de cette structure de torseur : sur
$\Gamma^+ \times \Gamma^+$ c'est celle de $\Gamma^+$ ; sur $\Gamma^- \times
\Gamma^-$, $\sigma\sigma'$ est l'élément de $\Gamma^+$ qui envoie $\sigma'$
sur $\sigma$ ; sur $\Gamma^+ \times \Gamma^-$ c'est l'action du torseur ; et
$\sigma v = (-v)\sigma$\footnote{Pour e), la page ajoute « si $\Gamma^+$ est
cyclique d'ordre $n$, ceci signifie $n \neq 1$ » ; il faut $n \geq 3$, car pour
$n = 2$, $\Gamma^+ \setminus {}_2\Gamma^+$ est vide. Pour $\Gamma^- \times
\Gamma^-$, la page écrit « l'application ct. $\to 1 \in \Gamma^+$ » ; le
produit n'est constant que sur la diagonale ($\sigma^2 = 1$). La lettre e)
surcharge un premier signe et vient avant d) sur la page.}.

Inversement, un groupe abélien $\Gamma^+$ et un torseur $\Gamma^-$ sous lui
définissent ainsi un groupe : on obtient une équivalence entre la catégorie des
torseurs sous des groupes abéliens et celle des extensions « diédrales »
scindées de $\mathbf{Z}/2$ par des groupes abéliens, pleinement fidèle, sur les
couples satisfaisant e), vers les groupes du début.

\textbf{Proposition} (page 39). Pour $n \geq 3$, la gerbe des formes de
$\mathbb{D}_n$ est équivalente à la gerbe des droites affines sur l'anneau
$\mathbf{Z}/n$ ; en particulier $\mathrm{Aut}(\mathbb{D}_n) \simeq
\mathrm{Aff}(1, \mathbf{Z}/n) \simeq (\mathbf{Z}/n)^\times \ltimes
\mathbf{Z}/n$\footnote{« Gerbe », écrit deux fois au-dessus de « cat. » biffé,
est une lecture incertaine.}.

\pagerange{40}{40}
\subsection*{Le groupe de l'icosaèdre sur $\mathbf{Z}$ (page 40)}

\textbf{Théorème} (page 40). Soit $S$ un schéma quelconque. Pour toute forme
$G$ de $\mathrm{PGL}_{2,S}$ et tout monomorphisme $\varphi :
\mathfrak{A}_{5,S} \hookrightarrow G$, posons $\alpha(\varphi) =
A(\varphi(\pi))$, $\pi$ le $5$-cycle standard. Alors le foncteur $(G, \varphi)
\mapsto \alpha(\varphi)$, du groupoïde des $(G, \varphi)$ vers la catégorie
discrète des $\alpha \in \Gamma(S, \mathcal{O}_S)$ tels que $\alpha^2 + \alpha
- 1 = 0$, est une équivalence. Autrement dit le groupoïde est \emph{rigide},
et le foncteur des classes d'isomorphisme est représentable par
\[
  M = \mathrm{Spec}\, \mathbf{Z}[T]/(T^2 + T - 1) .
\]
\textbf{Corollaire.} Si $u$ est un automorphisme intérieur de
$\mathfrak{A}_5$, $\alpha(\varphi \circ u) = \alpha(\varphi)$ ; s'il n'est
intérieur sur aucune fibre, $\alpha(\varphi \circ u) = \alpha'(\varphi) = -1 -
\alpha(\varphi) = -1/\alpha(\varphi)$.

En effet un automorphisme extérieur envoie $\pi$ sur un conjugué de $\pi^2$, et
$A(\pi^2) = S_2(\alpha) = \alpha^2 - 2 = -1 - \alpha$. C'est le théorème de la
page 10 débarrassé du facteur $\mathbf{Z}/2$\footnote{Un premier énoncé, en
termes de bijection sur les classes d'isomorphisme, est encadré et barré. La
note marginale « il suffit que \ldots{} soit trivial sur les fibres » vise sans
doute l'hypothèse de monomorphisme, qu'il suffit de vérifier fibre à fibre.}.

\pagerange{42}{56}
\section*{VI. Systèmes de Schwarz dans $\mathfrak{A}_4$, $\mathfrak{S}_4$, $\mathfrak{A}_5$ (pages 42 à 56)}

Un \emph{système de Schwarz} d'un groupe fini $H$ est un triple $(\rho, \rho',
\rho'') \in H^3$ de produit $\rho\rho'\rho'' = 1$ qui engendre $H$ ; son
\emph{type numérique} est le triple des ordres. C'est un homomorphisme
surjectif du groupe $\langle \ell_0, \ell_1, \ell_2 \mid \ell_0\ell_1\ell_2 =
1\rangle$ sur $H$, c'est-à-dire, géométriquement, un revêtement galoisien de
la sphère privée de trois points de groupe $H$ (section IX)\footnote{Le nom
renvoie aux triangles de Schwarz, dont les groupes de réflexions engendrent
les groupes finis de la sphère ; il l'écrit « Schwarz » à la page 1 et
« Schwartz » à partir de la page 41. L'interprétation par les revêtements est
de l'édition, mais elle est dans le plan de la page 1 (rubrique (1')).}.

\pagerange{42}{42}
\subsection*{Le tableau des types (page 42)}

La page 42 dresse les types numériques de systèmes de Schwarz, en
distinguant les deux classes de conjugaison d'éléments d'ordre $3$ de
$\mathfrak{A}_4$ (notées $3$, $3'$) et d'ordre $5$ de $\mathfrak{A}_5$ ($5$,
$5'$) :
\begin{itemize}
  \item $\mathfrak{A}_4$ (4 types) : $(3, 3', 2)$ et ses permutations
    circulaires (3 types), $(3, 3, 3)$ ;
  \item $\mathfrak{S}_4$ (9 types) : les six permutations de $(4, 3, 2)$
    (« standard ») et les trois de $(4, 4, 3)$ ;
  \item $\mathfrak{A}_5$ (32 types) : les six permutations de $(5, 3, 2)$
    (« standard $\alpha$ ») et de $(5', 3, 2)$ ; $(5, 5, 3)$ et $(5', 5', 3)$
    avec leurs permutations circulaires ; les six permutations de $(5, 5', 3)$
    et de $(5, 5', 2)$ ; $(5, 5, 5)$ ; $(5', 5', 5')$.
\end{itemize}
Total $45$, « ou 16 si on regarde modulo action de $\mathfrak{S}_5$ ». La
question posée en dessous : un système de Schwarz d'un groupe fini est-il
déterminé, à conjugaison près, par la classe de conjugaison de chacun de ses
termes ? « OK dans $\mathfrak{A}_4$, $\mathfrak{S}_4$, $\mathfrak{A}_5$ ---
regarder groupes diédraux, et $\mathfrak{S}_5$ ».

Les deux premières lignes sont exactes : $\mathfrak{A}_4$ a $96$ systèmes de
Schwarz, en $4$ classes sous $\mathrm{Aut}(\mathfrak{A}_4) \simeq
\mathfrak{S}_4$ ; $\mathfrak{S}_4$ en a $216$, en $9$ classes. La troisième
est incomplète : $\mathfrak{A}_5$ a $2280$ systèmes, en $38$ classes sous
$\mathfrak{A}_5$ et $19$ sous $\mathfrak{S}_5$, et non $32$ et $16$ ; il manque
le type $(3, 3, 5)$, avec ses permutations circulaires, en deux versions
$(3, 3, 5)$ et $(3, 3, 5')$, soit six classes\footnote{Ces nombres ont été
recalculés par énumération pour cette édition. Les $19$ classes sous
$\mathrm{Aut}(\mathfrak{A}_5)$ sont les $19$ « $T_2$-systèmes » de Philip Hall
(« The Eulerian functions of a group », 1936). La
réponse à la question de la page est bien « oui » pour les trois groupes,
$(3, 3, 5)$ compris : dans le recensement, chaque type, $5$ et $5'$
distingués, apparaît une seule fois. Les effectifs $4$, $9$, $32$, $45$ sont
dans la marge ; la lecture du premier triplet du crochet de la première ligne
est incertaine.}. L'omission se propage aux tables des pages 66 et 67
(section VII).

\pagerange{43}{44}
\subsection*{Formes de $\mathfrak{A}_4$ et de $\mathfrak{S}_4$ (pages 43 et 44)}

Soit $X$ une conique sur $S$ de caractéristique $0$ et $\mathbb{A} =
\underline{\mathrm{Aut}}(X)$. Les sous-groupes $\Gamma^+ \subset \Gamma$ de
$\mathbb{A}$ qui sont des formes de $\mathfrak{A}_4$ et de $\mathfrak{S}_4$ se
correspondent par $\Gamma = \mathrm{Norm}_{\mathbb{A}}(\Gamma^+)$, $\Gamma^+ =
[\Gamma, \Gamma]$. Comme $\mathrm{Aut}(\mathfrak{A}_4) \simeq
\mathrm{Aut}(\mathfrak{S}_4) \simeq \mathfrak{S}_4$, les catégories des formes
de $\mathfrak{A}_4$, des formes de $\mathfrak{S}_4$ et des ensembles à quatre
éléments sont équivalentes : à $E$ on associe $\mathfrak{S}_E$ et
$\mathfrak{A}_E$ ; à $\Gamma^+$ son groupe d'automorphismes et l'ensemble de
ses sous-groupes d'ordre $3$ ; à $\Gamma$ son groupe dérivé et l'ensemble de
ses sous-groupes d'ordre $6$. Donc sur $S$
\[
  \mathrm{Rev}_4(S) \simeq F(S, \mathfrak{S}_4) \simeq F(S, \mathfrak{A}_4)
  \simeq \mathrm{Tors}(\mathfrak{S}_{4,S}) \simeq \{\text{cubes orientés sur }
  S\} .
\]
Tout automorphisme de $\mathbb{A}$ étant intérieur et le normalisateur étant
$\Gamma$ lui-même, $\Gamma \xrightarrow{\sim} \underline{\mathrm{Aut}}(X,
\Gamma) = \underline{\mathrm{Aut}}(X, \Gamma^+)$ ; la catégorie des couples
$(X, \Gamma)$ est donc encore équivalente à celle des torseurs sous
$\mathfrak{S}_4$, des cubes orientés, des tétraèdres (non orientés). Dans ce
dictionnaire, $\Gamma$ est le groupe des automorphismes du tétraèdre ou des
rotations du cube, $\Gamma^+$ celui des automorphismes orientés du
tétraèdre\footnote{C'est le corollaire principal des pages 34 et 35, retrouvé.
La fin de la page 44 (le rôle des cubes « bi-orientés ») se lit mal ; « Re
NB » s'ouvre par « si on est en car. 0 (inutile ?) ».}.

On s'intéresse alors aux systèmes $(\mathbb{A}, \mathbb{H}, \rho, \rho',
\rho'')$ où $\rho, \rho', \rho''$ sont des sections de $\mathbb{A}$ de produit
$1$ qui engendrent un sous-groupe fini $\mathbb{H}$ de type $\mathfrak{A}_4$,
$\mathfrak{S}_4$ ou $\mathfrak{A}_5$. Un automorphisme d'un tel système
centralise $\mathbb{H}$, donc est trivial : la catégorie est \emph{rigide}, et
ses objets sont uniques à isomorphisme unique près une fois fixée leur classe.
Le but : trouver les systèmes de trois générateurs du groupe abstrait des
automorphismes orientés du tétraèdre, du cube, de l'icosaèdre\footnote{La page
44 est d'une prose rapide, en grande partie illisible entre les formules ;
« splittée » (deux fois) est une lecture incertaine.}.

\pagerange{45}{48}
\subsection*{L'action de $\mathfrak{S}_3$ ; systèmes canoniques d'un polyèdre (pages 45 à 48)}

Le premier invariant d'un système est son type numérique $(\nu, \nu',
\nu'')$, avec $2 \leq \nu, \nu', \nu'' \leq 3$, $4$ ou $5$ dans les cas I
($\mathfrak{A}_4$), II ($\mathfrak{S}_4$), III ($\mathfrak{A}_5$) ; le type
contient $3$, $4$ ou $5$ selon le cas, ce qui le reconnaît. Dans les cas I et
II le système est déterminé par son type ; dans le cas III il faut distinguer
les deux classes d'éléments d'ordre $5$ (page 45).

Sur l'ensemble $\mathrm{Tr}(H)$ des triples de produit $1$ opère
$\mathfrak{S}_3$ : l'orbite de $(\rho, \rho', \rho'')$ est formée de ses trois
permutations circulaires et des trois triples $(\rho''^{-1}, \rho'^{-1},
\rho^{-1})$, $(\rho'^{-1}, \rho^{-1}, \rho''^{-1})$, $(\rho^{-1}, \rho''^{-1},
\rho'^{-1})$. On veut savoir quels transformés d'un système lui sont
isomorphes, l'isomorphie étant celle que définit le groupe ambiant
$\mathbb{A}$ : conjugaison par $\mathfrak{S}_4$ (et non $\mathfrak{A}_4$) dans le
cas I, automorphismes intérieurs dans les cas II et III, avec en III une
classification plus grossière modulo $\mathfrak{S}_5$ (page 46).

Soient $\Pi$ un polyèdre combinatoire régulier orientable, $H =
\mathrm{Aut}^+(\Pi)$ et $\widetilde H = \mathrm{Aut}(\Pi)$. L'orientation
attache à chaque facette $\varphi \in S \amalg A \amalg F$ (sommets, arêtes,
faces) une rotation $\rho_\varphi \in H$, et à chaque repère (drapeau)
\emph{direct} $r = (s, a, f)$ le système
\[
  \Sigma_r = (\rho_f, \rho_s, \rho_a), \qquad \rho_f\rho_s\rho_a = 1, \quad
  \rho_a^2 = 1 ,
\]
le système \emph{canonique} ; les repères indirects donnent les systèmes
\emph{anticanoniques}, formés des inverses. $\Sigma_r$ ne dépend pas d'une
orientation choisie d'avance, puisque $r$ en définit une. Les $\Sigma_r$, $r$
direct, sont conjugués entre eux sous $H$ ; canoniques et anticanoniques le
sont sous $\widetilde H$, pas nécessairement sous $H$ --- il faut pour cela que
$\rho_s$ et $\rho_s^{-1}$ soient conjugués dans $H$, ce qui manque pour
l'icosaèdre. Enfin si $r_2 = g \cdot r_1$, $\Sigma_{r_2} = \mathrm{int}(g)
\Sigma_{r_1}$, donc $\Sigma_{r_1} = \Sigma_{r_2}$ si et seulement si $g$
centralise $H$ dans $\widetilde H$ ; pour l'octaèdre et l'icosaèdre, où
$\widetilde H = \mathbf{Z}/2 \times H$, c'est le cas de l'antipodie (pages 47 et
48)\footnote{Le début du « Théorème » de la page 45 sur le tétraèdre, et le haut
de la page 47 jusqu'à la fin de a), sont barrés de longs traits obliques ; la
page 46 est un feuillet inséré. Page 48, le facteur devant
$\mathrm{Aut}^+(\Pi)$ est lu « 3 » ; c'est $\mathbf{Z}/2 \times$. L'ordre des
deux premiers termes du système anticanonique varie (page 48 :
$(\rho_f^{-1}, \rho_s^{-1}, \rho_a)$ ; page 49 : $(\rho_s^{-1}, \rho_f^{-1},
\rho_a)$, qui est le transformé de $\Sigma_r$ par une transposition de
$\mathfrak{S}_3$) ; on ne fixe pas le choix. La page 47 est raturée en
plusieurs couches et les limites des passages biffés sont incertaines.}.

\pagerange{49}{52}
\subsection*{Le tétraèdre (pages 49 à 52)}

\textbf{Proposition} (pages 49 et 50). Soient $\Pi$ un tétraèdre orienté, $H =
\mathrm{Aut}^+(\Pi) \simeq \mathfrak{A}_4$, $\widetilde H = \mathrm{Aut}(\Pi)
\simeq \mathfrak{S}_4$.
\begin{itemize}
  \item[a)] $\varphi \mapsto \rho_\varphi$ est une bijection de $S \amalg F$
    sur les $8$ éléments d'ordre $3$ de $H$ ; si $s$ et $f$ sont opposés,
    $\rho_s\rho_f = 1$.
  \item[b)] $a \mapsto \rho_a$ est une application de degré $2$ des $6$
    arêtes sur les $3$ involutions de $H$, deux arêtes ayant même image si et
    seulement si elles sont opposées.
  \item[c)] $r \mapsto \Sigma_r$ envoie les repères sur les systèmes de type
    $(3, 3, 2)$ : son image est une orbite sous $\widetilde H$ (un torseur) et
    deux sous $H$.
  \item[d)] Les six transformés d'un système canonique par $\mathfrak{S}_3$
    sont de types distincts deux à deux, sauf pour la paire $\Sigma_r$,
    $\Sigma_{r'}$ ($r'$ le repère opposé), de même type et conjugués sous
    $\widetilde H$ mais non sous $H$.
  \item[e)] Si $s, s', s''$ forment une face $f$ et sont rangés dans l'ordre
    opposé à l'orientation qu'induit $\Pi$, $(\rho_s, \rho_{s'}, \rho_{s''})$
    est un système de type $(3, 3, 3)$, et de même pour trois faces ;
    chacun des deux paquets de $12$ forme une orbite (un torseur) sous $H$, et
    leur réunion un torseur sous $\widetilde H$.
  \item[f)] Ces $72 + 24 = 96$ systèmes sont tous les systèmes de Schwarz de
    $H$.
\end{itemize}
Les marges comptent : $24 = 12 + 12$ systèmes canoniques et anticanoniques,
$72 = 36 + 36$ avec leurs permutations circulaires\footnote{En b), la page
écrit « l'ens. ${}_3H$ des él. d'ordre 3 » ; ce sont les éléments d'ordre $2$.
La fin de la page 49, à partir de d), est très surchargée, et l'énoncé d) est
reconstitué à partir des fragments lisibles et de la suite. Les notes
marginales de la page 50 ne se lisent que par fragments. Le total $96$ est
confirmé par énumération.}.

\emph{Démonstration de f).} Les éléments $\neq 1$ de $H$ sont d'ordre $2$ ou
$3$, et les involutions forment, avec $1$, le sous-groupe $V$ d'indice $3$ ;
un système ne peut donc contenir deux involutions, et contient au moins deux
éléments d'ordre $3$, chacun un $\rho_s$ ou un $\rho_f$. Quatre cas, de $12$
systèmes chacun : un $\rho_f$ et un $\rho_s$ avec $s$ sur $f$ (si $s$ est
opposé à $f$, $\rho_s = \rho_f^{-1}$ et les deux n'engendrent pas $H$) : le cas
canonique ; la même chose dans l'autre orientation : l'anticanonique ; deux
sommets ; deux faces : les cas de e)\footnote{La page dit le sous-groupe des
involutions « d'indice 2 » ; $V$ est d'indice $3$ dans $\mathfrak{A}_4$, ce qui
ne change rien à l'argument. Les lettres finales de 3) et 4) sont écrites sur
d'autres, raturées.}.

\textbf{Corollaire} (page 51). Les types possibles pour $H \simeq \mathfrak{A}_4$
sont $(3, 3, 2)$, $(3, 2, 3)$, $(2, 3, 3)$ (canoniques et anticanoniques) et
$(3, 3, 3)$ ; la classe sous $\widetilde H = \mathrm{Aut}(H) \simeq
\mathfrak{S}_4$ est déterminée par le type. Un $H$ muni d'un système de Schwarz
est donc, $H = \mathrm{Aut}^+(\Pi)$, ou bien un repère $r$ de $\Pi$ avec
$\Sigma = \Sigma_r$, ou bien trois sommets distincts $(s, s', s'')$ --- ce qui
revient à un repère $\{s, \{s, s'\}, \{s, s', s''\}\}$ --- avec $\Sigma =
(\rho_s, \rho_{s'}, \rho_{s''})$ calculé dans l'orientation qui induit sur la
face $\{s, s', s''\}$ l'ordre opposé\footnote{La dernière ligne de la page 51 et
le haut de la page 52 sont surchargés et barrés ; on donne la formulation
retenue sur la page 52.}.

\pagerange{52}{56}
\subsection*{L'octaèdre (pages 52 à 56)}

\textbf{Proposition} (pages 52 et 53). Soient $\Pi$ un octaèdre combinatoire
orienté, $H = \mathrm{Aut}^+(\Pi) \simeq \mathfrak{S}_4$, $H^+ = [H, H] \simeq
\mathfrak{A}_4$, $R$ et $R^+$ les ensembles des repères et des repères directs.
\begin{itemize}
  \item[a)] Pour $r \in R$, $\Sigma_r$ est de type $(3, 4, 2)$ ; $\Sigma_r =
    \Sigma_{r'}$ si et seulement si $r$ et $r'$ sont antipodaux ; les $\Sigma_r$
    forment un torseur sous $H$, en bijection avec $R^+$ ($24$ systèmes, $144$
    avec les transformés par $\mathfrak{S}_3$).
  \item[b)] $s \mapsto \rho_s$ est une bijection des $6$ sommets sur les
    éléments d'ordre $4$, $f \mapsto \rho_f$ des $8$ faces sur les éléments
    d'ordre $3$, et $a \mapsto \rho_a$ fait des $12$ arêtes un revêtement de
    degré $2$ des $6$ involutions hors de $H^+$, deux arêtes ayant même image si
    et seulement si elles sont antipodales ; l'antipodie correspond au passage
    à l'inverse.
  \item[c)] Pour toute arête $\{s, t\}$, soit $f$ la face telle que $(s, \{s,
    t\}, f) \in R^+$ et $f'$ la face antipodale : $(\rho_s, \rho_t, \rho_{f'})$
    est un système de type $(4, 4, 3)$, et ces systèmes forment un torseur
    sous $H$ ($24$, $72$ avec les permutations circulaires).
  \item[d)] Il n'y en a pas d'autres.
\end{itemize}
\textbf{Corollaire} (page 54). Les types possibles sont les six permutations de
$(4, 3, 2)$ et les trois permutations circulaires de $(4, 4, 3)$, et le type
détermine le système à conjugaison près : un système est, d'une seule
manière, le transformé d'un $\Sigma_r$ ($r \in R^+$) par un $\sigma \in
\mathfrak{S}_3$, ou celui d'un $(\rho_s, \rho_t, \rho_{f'})$ par un $\sigma \in
\mathfrak{S}_3^+$\footnote{Le type « $(3, 4, 2)$ » ajouté en a) est surchargé et
de lecture incertaine ; il est juste. Les totaux $144 + 72 = 216$ sont ceux de
l'énumération. En c), l'énoncé porte $\rho_{f'}$ dans la formule, puis $\rho_f$
dans la discussion des transformés ; on suit la formule.}.

\emph{Démonstration de d)} (pages 54 à 56). Les ordres des éléments $\neq 1$
sont $2$, $3$, $4$. 1) Un système contient un élément d'ordre $4$ : il ne
contient pas deux involutions (sinon $H$ serait diédral), donc il en
contiendrait deux d'ordre $3$, qui sont dans $H^+$ --- impossible, ils
n'engendreraient pas $H$. 2) Il contient un élément d'ordre $3$ : sinon, ne
contenant pas deux involutions, il contiendrait deux éléments d'ordre $4$,
$\rho_s$ et $\rho_t$ quitte à agir par $\mathfrak{S}_3$, avec $t \neq s$ et
$t$ non antipodal à $s$, donc $\{s, t\}$ est une arête, et l'on est dans c), où
le troisième terme est d'ordre $3$. 3) Restent $\rho = \rho_s$, $\rho' =
\rho_f$ : si $s$ est sur $f$, c'est le cas a) ; sinon la seule autre position
relative d'un élément d'ordre $4$ et d'un d'ordre $3$ qui engendrent est celle
de c)\footnote{L'ordre de la démonstration sur les pages est 54 (début), 56
(points 1 et 2, et le début de 3), 55 (la fin) ; on la donne dans l'ordre
logique. Un passage de cinq lignes est encadré et biffé page 56, un autre
page 55. Le mot « absurde » de la page 55 est une lecture incertaine.}.

\pagerange{58}{73}
\section*{VII. Hyperpolygônes réguliers (pages 58 à 73)}

Un « hyperpolygône » régulier est la donnée, dans un plan vectoriel $V$ sur un
corps, de deux pseudo-réflexions $\tau_0, \tau_1$ --- des automorphismes
ayant une droite de points fixes, de valeur propre non triviale $\lambda_i$ ---
qui engendrent un groupe et satisfont des conditions de non-dégénérescence ;
quand $\lambda_0 = \lambda_1 = -1$ ce sont deux réflexions, et l'on retrouve
les polygones réguliers ordinaires. Ce sont, sous le nom que leur a donné
Coxeter, les \emph{polygones réguliers complexes}, et leurs groupes sont les
groupes de réflexions complexes de rang $2$ de Shephard et Todd\footnote{Le
rapprochement est de l'édition. G. C. Shephard (1952) et H. S. M. Coxeter
(\emph{Regular Complex Polytopes}, 1974) notent $p\{q\}r$ ces polygones et
$p[q]r$ leurs groupes ; Shephard et Todd (1954) ont classé les groupes finis
engendrés par des pseudo-réflexions. Rien n'indique que Grothendieck les ait
connus ; il ne les cite pas.}. Dans la section VI, les sommets et les faces
d'un polyèdre donnaient les systèmes de Schwarz ; ici, ce sont les images
$\tau_0^C, \tau_1^C$ dans $\mathrm{Aut}(C) \simeq \mathrm{PGL}_2$, $C =
\check{\mathbb{P}}(V)$, complétées par l'inverse de leur produit, qui forment
un tel système.

\pagerange{58}{60}
\subsection*{Le cas d'un diviseur invariant ; la borne $\nu_i \leq 5$ (pages 58 à 60)}

Sur un corps algébriquement clos de caractéristique $0$, avec $G_0 = \langle
\tau_0, \tau_1\rangle$ fini, on écrit $\tau_i = 1 + a_i' \otimes a_i$ ($a_i \in
V$, $a_i' \in V^\vee$), $\mu_i = \langle a_i, a_i'\rangle = \lambda_i - 1$, et
\[
  \beta = \langle a_0, a_1'\rangle\langle a_1, a_0'\rangle
\]
l'invariant du couple. Les $\lambda_i$ sont des racines de l'unité d'ordres
$\nu_i \geq 2$ ; les conditions de non-dégénérescence sont $\langle a_0,
a_1'\rangle \neq 0$ et $a_0, a_1$ linéairement indépendants\footnote{La page 58
passe d'abord par un espace $\mathcal{E} \simeq V \oplus k\pi$, sous la
condition $\langle \pi, \pi'\rangle = \mu_0\mu_1 - \beta \neq 0$, puis se borne
à $V$. C'est le cas $n = 1$ des « hyperpolyèdres » des dossiers 70 et 72.}.

I) \emph{Si la représentation sur $C$ admet un diviseur invariant de degré
$2$}, alors $\tau_0, \tau_1$ sont d'ordre $2$ : c'est le cas des polygones
réguliers ordinaires. Le diviseur n'est pas un point double (une droite
invariante aurait un supplémentaire invariant et $G_0$ serait commutatif) ;
c'est donc une paire de droites $V = D \oplus D'$, qu'aucun $\tau_i$ ne
laisse invariantes toutes deux. Si un $\tau_i$ les fixait, on aurait, quitte à
échanger, $D = ka_i$, $D' = \mathrm{Ker}\, a_i'$, et pour $j \neq i$, $\tau_j$
les échangerait ; en écrivant $\tau_j D \subset D'$ on trouve $\mu_i + \beta =
0$, et en écrivant $\tau_j^2 D \subset D$ (avec $\tau_j^2 = 1 + (2 + \mu_j)\,
a_j' \otimes a_j$), $\mu_i + (2 + \mu_j)\beta = 0$, d'où $(1 + \mu_j)\beta = 0$,
$\beta \neq 0$, donc $\lambda_j = 0$ : absurde. Donc $\tau_0$ et $\tau_1$
échangent $D$ et $D'$ ; dans une base $e_0 \in D$, $e_1 = \tau_0 e_0$,
\[
  \tau_0 = \begin{pmatrix} 0 & x \\ 1 & 0 \end{pmatrix}, \qquad
  \tau_1 = \begin{pmatrix} 0 & z \\ y & 0 \end{pmatrix},
\]
de polynômes caractéristiques $\lambda^2 - x$ et $\lambda^2 - yz$ ; comme $1$
est valeur propre, l'autre est $-1$\footnote{La page note $\xi, \zeta, \eta$
les coefficients et écrit les équations aux valeurs propres $\lambda^2 + \xi =
0$, $\lambda^2 + \zeta\eta = 0$ ; avec ce signe la valeur propre $1$ ne
s'obtient pas. Le signe est corrigé ; la conclusion est celle de la page.}.

II) \emph{Sinon}, l'image de $G_0$ dans $\mathrm{Aut}(C) \simeq
\mathrm{PGL}_2(k)$ est l'un des groupes $\mathfrak{A}_4$, $\mathfrak{S}_4$,
$\mathfrak{A}_5$ (section V : les cycliques et diédraux ont un diviseur
invariant de degré $2$), et $G_0 \subset \widetilde{G}_0 \cdot k^\times$, où
$\widetilde{G}_0$ est l'image réciproque dans $\mathrm{SL}(V)$. Le groupe
$Z_i$ engendré par $\tau_i$ rencontre les homothéties trivialement, donc
s'injecte dans $\mathrm{Aut}(C)$, et
\[
  2 \leq \nu_0, \nu_1 \leq 5 ;
\]
de même pour toute pseudo-réflexion de $G_0$ (page 60)\footnote{La lettre du
groupe engendré par $\tau_i$ a la forme d'un $3$ bouclé ; la transcription la
rend $Z_i$. La page 60 s'arrête sur « compris entre 2 et 5 \ldots ».}.

\pagerange{61}{62}
\subsection*{L'équation de $\xi$ (pages 61 et 62)}

On travaille sur $\mathbf{C}$, avec $\lambda_i = e^{i\theta_i}$, $0 < \theta_i
\leq \pi$, et l'on écrit $\beta = \xi\mu_0\mu_1$, $\xi$ réel\footnote{Le début
de la page 61 est surchargé ; on lit $\beta = c\,(\lambda_0 - 1)(\lambda_1 -
1)/(\kappa_1 - 1)$ et $\xi = c/(\kappa_1 - 1) \in \mathbf{R}$, sans que la
provenance de $c$ se lise. Que $\xi$ soit réel est affirmé, non démontré
ici.}. Le produit $u = \tau_0\tau_1$ a pour trace $\lambda_0 + \lambda_1 +
\beta$ et pour déterminant $\lambda_0\lambda_1$, et l'invariant de son image
dans $\mathrm{PGL}_2$ est, par la page 27,
\[
  \kappa = \frac{(\mathrm{Tr}\, u)^2}{\det u} - 2 \in \{-2, -1, 0, c_5, c_5'\}
\]
selon que $\tau_0^C\tau_1^C$ est d'ordre $2$, $3$, $4$ ou $5$ (deux classes).
Cela s'écrit
\[
  \frac{\mu_0^2\mu_1^2}{\lambda_0\lambda_1}\Bigl[\xi^2 + 2\,
  \frac{\lambda_0 + \lambda_1}{\mu_0\mu_1}\,\xi + \frac{\lambda_0^2 +
  \lambda_1^2}{\mu_0^2\mu_1^2}\Bigr] = \kappa ,
\]
et, avec $|\lambda - 1| = 2\sin\frac\theta2$, $\mu_0^2\mu_1^2/\lambda_0\lambda_1
= 16 s_0^2 s_1^2$ où $s_i = \sin\frac{\theta_i}{2}$, l'équation devient
\[
  (\mathrm E) \qquad \xi^2 - \frac{\cos\frac{\theta_1 - \theta_0}{2}}{s_0 s_1}\,
  \xi + \frac{2\cos(\theta_1 - \theta_0) - \kappa}{16\, s_0^2 s_1^2} = 0 ,
\]
de discriminant
\[
  \Delta(\kappa) = \frac{\kappa + 2}{4 s_0^2 s_1^2} \geq 0, \qquad \kappa + 2 =
  4\cos^2\frac\Theta2 \quad (\kappa = 2\cos\Theta) ,
\]
de racines
\[
  \xi_\pm = \frac{\cos\frac{\theta_1 - \theta_0}{2} \pm
  \cos\frac{\Theta}{2}}{2 s_0 s_1} .
\]
Ce sont les formules des pages 61, 62 et 70\footnote{La page écrit le terme
constant avec $\frac18$ au lieu de $\frac1{16}$, ici et page 62 ; son propre
calcul du discriminant, qui donne le $\Delta$ encadré, demande $\frac1{16}$. Le
facteur devant le crochet, écrit dans la marge, est lu avec doute
$(\lambda_0 - 1)(\lambda_1 - 1)/\lambda_0\lambda_1$ ; c'est son carré. Une
colonne de calculs en travers de la marge ne se lit qu'en partie.}.
Le discriminant ne s'annule que pour $\kappa = -2$, c'est-à-dire
$\tau_0^C\tau_1^C$ d'ordre $2$, où la racine est double, $\xi = \cos\frac{\theta_1
- \theta_0}{2}/(2s_0s_1)$. En prenant $\lambda_i = e^{2i\pi/\nu_i}$, $2 \leq
\nu_i \leq 5$, sans $\nu_0 = \nu_1 = 2$, il y a $15$ possibilités pour
$\kappa = -2$, et a priori $8 \cdot 15 = 120$ pour les quatre autres valeurs de
$\kappa$, deux racines chacune --- « mais il faut en exclure un bon
paquet ! »\footnote{Le « 16 » repassé, barré, avec « $= 15$ » au-dessus, est
lu avec doute ; $15 = 4 \cdot 4 - 1$ est le compte juste.}.

Deux principes d'élimination (page 62) : a) $0 < \xi < 1$, qui exprime que la
forme hermitienne invariante est définie (« ou $=$ ») ; b) $\xi$ doit encore
satisfaire les équations obtenues en remplaçant $\theta_0$ par $i\theta_0$
et $\theta_1$ par $j\theta_1$ ($i \not\equiv 0 \bmod \nu_0$, $j \not\equiv 0
\bmod \nu_1$), pour une valeur convenable de $\kappa$ dans la même liste ---
puisque les puissances $\tau_0^i, \tau_1^j$ engendrent le même groupe.

\pagerange{62}{65}
\subsection*{Situation projective et situation hyperpolygonale (pages 62 à 65)}

L'analyse des systèmes de Schwarz (section VI) montre que la \emph{situation
projective} --- le couple $(\tau_0^C, \tau_1^C)$ dans $C$, à isomorphisme
près --- est déterminée par $\kappa_0, \kappa_1, \kappa$, les classes de
conjugaison de $\tau_0^C$, $\tau_1^C$ et de leur produit dans $\mathrm{Aut}(X)$,
ou, ce qui revient au même, dans le groupe pythagoricien $H$ qu'ils
engendrent, modulo $\mathrm{Aut}(H)$ dans le cas tétraédral et
$\mathrm{Int}(H)$ sinon. Une situation projective donnée se relève en quatre
situations hyperpolygonales : il faut choisir l'un des deux points fixes de
$\tau_0^C$ et de $\tau_1^C$ ; c'est remplacer $\tau_i$ par
$\lambda_i^{-1}\tau_i$, ce qui change $\lambda_i$ en $\lambda_i^{-1}$ sans
changer $\kappa_0$, $\kappa_1$, $\kappa$ (page 63). Le dictionnaire de la
page 64 :
\[
  \{\text{données } (\lambda_0, \lambda_1, \beta)\} \leftrightarrow
  \{\text{couples } (\tau_0, \tau_1) \text{ épinglés}\} \leftrightarrow
  \{(\kappa_0, \kappa_1, \kappa) \text{ admissible}, + a_0, a_1\} ,
\]
où $a_i$ est un point fixe de $\tau_i^C$, et $\lambda_i$ l'une des deux racines
de $\lambda^2 - \kappa_i\lambda + 1 = 0$ ; le choix de $\lambda_i$ équivaut à
celui de $a_i$ si $\nu_i \neq 2$, et ne dit rien si $\nu_i = 2$
($\lambda_i = -1$)\footnote{La page 64 écrit $\lambda_i^2 + \kappa_i\lambda_i +
1 = 0$ ; avec $\kappa_i = \lambda_i + \lambda_i^{-1}$ (page 61) le signe est
$-$. La moitié inférieure de la page porte, tête-bêche et barré, un brouillon
antérieur (trois faisceaux de droites, des formules en $\rho$, $\eta$,
$\lambda_0 = \varphi_0^2$) dont rien n'est repris.}.

\textbf{En résumé} (page 65). a) Pour $(\lambda_0, \lambda_1, \kappa)$
admissible avec $\nu_0, \nu_1 \neq 2$, il y a un et un seul hyperpolygône à
isomorphisme près, c'est-à-dire un seul $\xi$ : si $\kappa = -2$ c'est la
racine double, sinon une seule des deux racines de (E) convient --- « Comment
la distinguer ? Si on travaille dans $\mathbf{C}$, on peut songer à la
condition $0 < \xi < 1$ ». b) Si $\nu_0$ ou $\nu_1$ vaut $2$, il y a \emph{deux}
hyperpolygônes non isomorphes, correspondant aux deux racines distinctes de
(E) --- on sait alors que $\kappa \neq -2$\footnote{Dans b), le signe du
crochet final se lit $=$ ou $\neq$ ; $\neq$ est retenu, la phrase concluant à
deux racines distinctes.}. Le type numérique $(\nu_0, \nu_1, \nu)$, une fois
normalisés $\lambda_0 = e^{2i\pi/\nu_0}$, $\lambda_1 = e^{2i\pi/\nu_1}$, ne
détermine donc pas l'hyperpolygône, pour deux raisons : un $\nu_i = 2$ (page
66, liste 1) ; ou $\nu = 5$, où l'on peut choisir $\kappa = c_5$ ou $c_5'$
(liste 2), ce qu'on note $5$ et $5'$.

\pagerange{66}{68}
\subsection*{Le recensement (pages 66 à 68)}

Grothendieck dresse la table des types, en soulignant ceux qui donnent deux
hyperpolygônes et en portant en indice le nombre $\varphi(\nu_0)\varphi(\nu_1)$
de normalisations :
\begin{itemize}
  \item cas tétraédraux ($6$) : $(3,3,2)$, $\underline{(3,2,3)}$,
    $\underline{(2,3,3)}$, $(3,3,3)$ ;
  \item cas octaédraux ($13$) : $(4,3,2)$, $\underline{(3,2,4)}$,
    $\underline{(2,4,3)}$, $\underline{(2,3,4)}$, $(3,4,2)$,
    $\underline{(4,2,3)}$, $(4,4,3)$, $(4,3,4)$, $(3,4,4)$ ;
  \item cas icosaédraux ($25$) : $(5,3,2)$, $\underline{(3,2,5)}$,
    $\underline{(2,5,3)}$, $\underline{(2,3,5)}$, $(3,5,2)$,
    $\underline{(5,2,3)}$ ; $\underline{(3,2,5')}$, $\underline{(2,3,5')}$ ;
    $(5,5,3)$, $(5,3,5)$, $(3,5,5)$ ; $(5,3,5')$, $(3,5,5')$ ;
    $\underline{(2,5,5')}$, $\underline{(5,2,5')}$ ; $(5,5,5)$.
\end{itemize}
Au total $44$ cas, dont $14$ où le type numérique ne détermine pas le type
d'isomorphie. Il vérifie que le total des indices redonne $4 \cdot 45 = 180$ :
chacun des $45$ systèmes de Schwarz $(\kappa_0, \kappa_1, \kappa)$ donne $2
\times 2$ possibilités. Pour calculer les $\xi$ dans les $180$ cas, il
faudrait vérifier comment $\xi$ change quand on remplace $\tau_0, \tau_1$ par
$\tau_0^{\ell_0}, \tau_1^{\ell_1}$, ce qui ramènerait aux $45$ cas normalisés
--- « qu'il faudrait expliciter \ldots »\footnote{Deux réserves. La ligne
$(5,5,3)$, $(5,3,5)$, $(3,5,5)$ compte trois cas et porte « $+4$ » ; avec $3$, le
total icosaédral serait $24$. Surtout, la table hérite de l'omission de la
page 42 : il y manque $(3,3,5)$ et $(3,3,5')$, deux pseudo-réflexions d'ordre
$3$ dont le produit est d'ordre $5$ dans $\mathrm{PGL}_2$. Ils existent : le
groupe de Shephard--Todd $3[5]3$ ($G_{20}$, d'ordre $360$, d'image
$\mathfrak{A}_5$) en donne. Les nombres en indice sont écrits sous les triples ;
dans $(5,3,5')$ et $(3,5,5')$ les premiers chiffres sont surchargés. La note de
la page 67 se perd dans des surcharges.}.

\pagerange{69}{73}
\subsection*{Le test (pages 69 à 73)}

Pour décider laquelle des deux racines de (E) convient, Grothendieck teste la
condition $0 < \xi < 1$. Dans les sept cas sans $2$ --- $(3,3,3)$ ; $(4,4,3)$,
$(4,3,4)$, $(3,4,4)$ ; $(5,5,3)$, $(5,3,5)$, $(3,5,5)$ ; $(3,5,5')$, $(5,3,5')$ ;
$(5,5,5)$ --- où $\frac{2\pi}{5} \leq \theta_0, \theta_1, \Theta \leq
\frac{2\pi}{3}$, les termes constants de (E) sont positifs, et il s'agit de
vérifier que la grande racine est $\geq 1$ (page 70), soit
\[
  \cos\frac{\theta_1 - \theta_0}{2} + \cos\frac\Theta2 \geq 2\sin\frac{\theta_0}{2}
  \sin\frac{\theta_1}{2} .
\]
Il le vérifie à la main : $(3,3,3)$, $\frac32 \geq \frac32$ (égalité) ;
$(4,4,3)$, $\frac32 \geq 1$ ; $(4,3,4)$ et $(3,4,4)$, $\sqrt{2 + \sqrt3} + \sqrt2
\geq \sqrt6$ ; $(5,5,3)$, $\frac32 \geq \frac12(2 - c_5)$ ; $(5,3,5)$ et $(3,5,5)$
par un calcul dans $\mathbf{Q}(\sqrt5)$ ; $(5,5,5)$. Pour $(5,3,5')$ et
$(3,5,5')$ il conclut « pas vrai ! », et ce sont les « seuls cas
indéterminés » de la page 73 : « vérifier s'il existe bien --- sinon, il faut
trouver comment les rendre éliminables ».

Les conclusions « OK » sont justes, et l'inégalité vaut aussi dans les deux
cas laissés ouverts : pour $(5,3,5')$, $\theta_0 = \frac{2\pi}{5}$, $\theta_1 =
\frac{2\pi}{3}$, $\Theta = \frac{4\pi}{5}$, le premier membre vaut $\cos\frac{2\pi}{15}
+ \cos\frac{2\pi}{5} \approx 1{,}223$ et le second $2\sin\frac\pi5\sin\frac\pi3
\approx 1{,}018$ ; les racines de (E) sont environ $0{,}594$ et $1{,}201$, et
une seule est dans $]0, 1[$\footnote{Valeurs calculées pour cette édition, ainsi
que, pour les autres cas, les racines $(\xi_-, \xi_+)$ : $(3,3,3)$ $(\frac13,
1)$ ; $(4,4,3)$ $(0{,}5 \,;\, 1{,}5)$ ; $(4,3,4)$ $(0{,}211 \,;\, 1{,}366)$ ; $(5,5,3)$
$(0{,}724 \,;\, 2{,}171)$ ; $(5,3,5)$ $(0{,}103 \,;\, 1{,}692)$ ; $(5,5,5)$ $(0{,}276 \,;\, 2{,}618)$. Dans les pages 71 et 72, les valeurs trigonométriques confondent
$\frac\pi5$ et $\frac{2\pi}5$ (« $\cos\frac\pi5 = \frac12\alpha$ »,
« $\sin^2\frac\pi5 = \frac14(3 + \alpha)$ », qui valent pour $\frac{2\pi}5$) ;
la page 71, qui fait le calcul de $(5,5,3)$ par les racines, est barrée en
entier. Le « pas vrai ! » de la page 72 vient d'une substitution
« $\cos\frac{2\pi}{5} = -\frac12$ » qui est fausse. Pour le type omis
$(3,3,5)$ la grande racine vaut environ $1{,}206$ ; pour $(3,3,5')$ les deux
racines, environ $0{,}461$ et $0{,}873$, sont dans $]0, 1[$, et le test ne
trancherait pas.}. Le cas $(3,3,3)$ est le cas limite $\xi_+ = 1$, que la
condition « ou $=$ » de la page 62 prévoyait.

\pagerange{75}{80}
\section*{VIII. Les $2$-polyèdres réguliers finis en caractéristique $0$ (pages 75 à 80)}

On revient aux polyèdres de la section II : trois réflexions $\tau_0, \tau_1,
\tau_2$ d'un espace de rang $3$, et l'on cherche quand $G_0 = \langle \tau_0,
\tau_1, \tau_2\rangle$ est fini.

\pagerange{75}{75}
\subsection*{Les invariants (page 75)}

Les invariants sont $\beta_0 = \alpha_0 + 2$, $\beta_1 = \alpha_1 + 2$, le
discriminant $\rho = 8 - 2(\beta_0 + \beta_1) = -2(\alpha_0 + \alpha_1)$,
$\tilde\rho = -(\alpha_0 + \alpha_1)$, et $\sigma = 4 - \beta_1 = 2 - \alpha_1$,
$\sigma' = 2 - \alpha_0$ (les notations du dossier 72). La conique $Q_C = Y$
du plan $C$ est lisse si et seulement si $\beta_0\beta_1\tilde\rho$ est
inversible, c'est-à-dire si en tout point $\alpha_0 \neq -2$, $\alpha_1 \neq -2$,
$\alpha_0 + \alpha_1 \neq 0$ --- la condition du théorème de la page 4 ; la
quadrique $Q_X$ de l'espace dual est lisse si et seulement si
$\beta_0\beta_1\sigma$ est inversible\footnote{La page conclut pour $Q_X$
« $\alpha_0 \neq -2$, $\alpha_1 \neq \pm 2$ », sans symétrie entre les indices ;
on le rapporte tel quel.}. $\rho_f = \tau_0\tau_1$ a un ordre qui dépend de
$\beta_0$ de la manière connue (section IV), $\rho_s = \tau_1\tau_2$ de même avec
$\beta_1$.

Comme $\mathrm{Aut}(C, Q_C) \xrightarrow{\sim} \mathrm{Aut}(Y)$, la catégorie des
$2$-polyèdres réguliers épinglés sur $S$ est équivalente à celle des couples
$(Y, \varphi^Y)$ d'une conique $Y$ et d'un homomorphisme $\varphi^Y : G_2 \to
\mathrm{Aut}(Y)$ tel que : a) $\rho_s^2 \neq 1$ et $\rho_f^2 \neq 1$ en toute
fibre (de sorte que les $\tau_i$ sont des réflexions) ; b) les points $h_0, h_1,
h_2$ de $C$ définis par les $\tau_i$ ne sont alignés en aucun point. Si en
aucun point on n'a $\alpha_1 = -2$, $\alpha_0 = 2$, b) équivaut à b') la
représentation sur $C$ n'a de droite fixe en aucun point ; et si $Y$ est lisse,
b') équivaut à b'') $Y$ n'a, sur aucune fibre, de diviseur de degré $2$
stable --- la non-dégénérescence de la section III\footnote{Le second terme de
l'idéal « $(\beta_1, \cdot)\mathcal{O}_S$ » est illisible ; la note marginale,
en biais et en partie biffée, discute la caractéristique $2$.}.

\pagerange{76}{77}
\subsection*{Le théorème : neuf polyèdres (pages 76 et 77)}

Sur un corps $k$, soient $\zeta_i$ des racines de $\zeta_i^2 - \alpha_i\zeta_i +
1 = 0$ dans $\bar k$ et $\nu_i$ l'ordre de $\zeta_i$.

\textbf{Proposition} (page 76). Si $G_0$ est fini, $\beta_0, \beta_1$ sont
algébriques sur le corps premier ; la réciproque vaut en caractéristique $p >
0$. De plus : a) si $\nu_i \neq 1, 2$ (c'est-à-dire $\alpha_i \neq \pm 2$),
$\nu_i$ est l'ordre de $\rho_i$ ; b) si $\alpha_i = -2$, $\rho_i$ est d'ordre
$2p$ ; c) si $\alpha_i = 2 \neq -2$, $\rho_i$ est d'ordre $p$ --- avec $p =
\infty$ en caractéristique nulle\footnote{La note qui suit « NB » et b) ne se
lisent que partiellement ; on donne la forme que leur donnent c) et la
section IV.}.

\textbf{Théorème} (page 76). Supposons $k$ de caractéristique $0$. $G_0$ est
fini si et seulement si l'on est dans l'un des six cas :
\begin{itemize}
  \item[I)] $\nu_0 = \nu_1 = 3$, $\alpha_0 = \alpha_1 = -1$ : le tétraèdre ;
  \item[II)] $(\nu_0, \nu_1) = (3, 4)$, $\alpha_0 = -1$, $\alpha_1 = 0$ :
    l'octaèdre ;
  \item[II')] $(\nu_0, \nu_1) = (4, 3)$ : le cube ;
  \item[III)] $(\nu_0, \nu_1) = (3, 5)$, $\alpha_0 = -1$, $\alpha_1^2 + \alpha_1
    - 1 = 0$ : l'icosaèdre ;
  \item[III')] $(\nu_0, \nu_1) = (5, 3)$ : le dodécaèdre ;
  \item[IV)] $\nu_0 = \nu_1 = 5$ et $\alpha_0 \neq \alpha_1$, solutions de
    $\alpha^2 + \alpha - 1 = 0$ : les cas « pentagonaux ».
\end{itemize}
Les cas I, II, II' donnent chacun une structure, définie sur $\mathbf{Q}$ ; les
trois autres en donnent chacune deux, qui existent si et seulement si $k$
contient une racine de $\alpha^2 + \alpha - 1$, c'est-à-dire $\sqrt5$. « Il y a
donc en tout $3 + 3 \times 2 = 9$ types d'isomorphie de \emph{polyèdres réguliers
finis} sur corps $k$ de caractéristique $0$, contenant $\alpha$ »
(page 77).

Ce sont les cinq solides de Platon et les quatre polyèdres de
Kepler--Poinsot\footnote{L'identification est de l'édition ; le dossier 70
(pages 64 à 71) la fait déjà, dans sa table des solides « et des
pentagrammiques ». Avec $\rho_f$ d'ordre $\nu_0$ (la face) et $\rho_s$ d'ordre
$\nu_1$ (la figure de sommet), l'invariant $c_5'$ correspond au pentagramme
$\frac52$. Que les polyèdres réguliers étoilés soient exactement quatre est le
théorème de Cauchy (1813), après Poinsot (1809). La page 77 dit « le corps des
racines cyclotomiques d'ordre 5 » ; c'est son sous-corps réel
$\mathbf{Q}(\sqrt5)$ qui suffit.} : en III, $\alpha_1 = c_5$ donne l'icosaèdre
$\{3, 5\}$ et $\alpha_1 = c_5'$ le grand icosaèdre $\{3, \frac52\}$ ; en III',
le dodécaèdre $\{5, 3\}$ et le grand dodécaèdre étoilé $\{\frac52, 3\}$ ; en
IV, $(c_5, c_5')$ et $(c_5', c_5)$ sont le grand dodécaèdre $\{5, \frac52\}$ et
le petit dodécaèdre étoilé $\{\frac52, 5\}$, duaux l'un de l'autre.

\pagerange{77}{80}
\subsection*{La démonstration (pages 77 à 80)}

On se ramène à $k = \mathbf{C}$.

I) Si $\nu_0, \nu_1$ sont finis, $\alpha_i \neq \pm 2$, et les $\alpha_i$,
racines d'équations semi-cyclotomiques d'indices $\geq 3$, sont réels : la
situation provient de $\mathbf{R}$. $G_0$ étant fini, il existe une forme
quadratique invariante définie positive, d'où $\beta\rho \neq 0$ et la
lissité. L'image $\underline{G}_0$ de $G_0$ dans $\mathrm{PGL}_2(\mathbf{C})$
est un groupe fini d'automorphismes de $Y \simeq \mathbb{P}^1$ ; par Hurwitz,
$Y \to Y/\underline{G}_0 \simeq \mathbb{P}^1$ est ramifié en trois points
exactement (avec deux points, un diviseur de degré $2$ serait stable,
contrairement à b'')), d'indices $(2, 2, n)$ --- exclu pour la même raison ---
ou $(2, 3, 3)$, $(2, 3, 4)$, $(2, 3, 5)$\footnote{La page attribue $(2, 3, 3)$
au tétraèdre. Mais le $G_0$ du tétraèdre est $\mathfrak{S}_4$ (le groupe de
Coxeter $[3, 3]$, cf. page 8), et son image dans $\mathrm{PGL}_2 \simeq
\mathrm{SO}(3)$ par $g \mapsto \det(g)\, g$ est encore $\mathfrak{S}_4$ : les
rotations-réflexions d'ordre $4$ du tétraèdre deviennent des rotations d'ordre
$4$. La signature est donc $(2, 3, 4)$ pour le tétraèdre comme pour
l'octaèdre, et $(2, 3, 3)$ ne se présente pas, $\mathfrak{A}_4$ n'étant pas
engendré par des involutions. La conclusion --- trois familles, tétraédrale,
octaédrale, icosaédrale --- n'en est pas changée, $(\nu_0, \nu_1)$ séparant
les deux premières. La note marginale précise que $G_0 \to \underline{G}_0$ a un
noyau d'ordre $1$ ou $2$.}.

II) Les cas I, II, III se réalisent : pour III on prend $\alpha = 2\cos\frac{2\pi}5$,
et l'on constate sur les formules explicites que l'on obtient des réalisations
géométriques strictes et fidèles du tétraèdre, de l'octaèdre et de l'icosaèdre.
Leurs groupes $G_0$ sont $\mathfrak{S}_4$, $\mathfrak{S}_4 \times \mathbf{Z}/2$,
$\mathfrak{A}_5 \times \mathbf{Z}/2$, d'images $\mathfrak{S}_4$, $\mathfrak{S}_4$,
$\mathfrak{A}_5$ dans $\mathrm{PGL}_2$ --- les « groupes pythagoriciens
standard »\footnote{La page lit le premier groupe « $\mathfrak{A}_4$ » avec doute
et donne les images $\mathfrak{A}_4$, $\mathfrak{S}_4$, $\mathfrak{A}_5$ ; voir la
note précédente et l'exemple A) de la page 8, qui donne $\mathfrak{S}_4$.} ; II'
et III' s'en déduisent par dualité, et IV par la construction des pentagones
associés à l'icosaèdre.

III) Dans les cas I, II, III, la construction donne des cartes cellulaires
régulières de la sphère simplement connexes : le groupe des rotations $G_0^+$
(${}\simeq \mathfrak{A}_4, \mathfrak{S}_4, \mathfrak{A}_5$) est présenté par
\[
  \rho_s^{\nu_0} = \rho_f^{\nu_1} = (\rho_s\rho_f)^2 = 1 ,
\]
et la théorie des groupes fondamentaux de la sphère privée de trois points
montre qu'en général $G_0$ est un quotient de l'un de ces groupes ; les ordres
des éléments étant $\leq 5$, $3 \leq \nu_0, \nu_1 \leq 5$. D'autre part la forme
quadratique invariante, de matrice dans la base $a_0, a_1, a_2$
\[
  \begin{pmatrix}
    2 & -\beta_0 & 0 \\
    -\beta_0 & 2\beta_0 & -\beta_0\beta_1 \\
    0 & -\beta_0\beta_1 & 2\beta_0\beta_1
  \end{pmatrix},
\]
de déterminant $2\beta_0^2\beta_1(4 - \beta_0 - \beta_1)$, est définie
positive si et seulement si $-2 < \alpha_0, \alpha_1 < 2$ et $\alpha_0 +
\alpha_1 < 0$ ; la première condition est automatique pour $\nu_i \geq 3$.
Quitte à dualiser, $3 \leq \nu_0 \leq \nu_1 \leq 5$. Restent, au-delà de I, II,
III, les couples $(4, 4)$, $(4, 5)$, $(5, 5)$. $(4, 4)$ donne $\alpha_0 +
\alpha_1 = 0$ : exclu. $(4, 5)$ avec $\alpha_1 = c_5 > 0$ est exclu ; $(5, 5)$
avec $\alpha_0 = c_5$ force $\alpha_1 = c_5'$ ($\alpha_0 + \alpha_1 = -1$),
c'est IV. « On a gagné ! » (page 80).

Il manque à cette fin d'argument trois couples, que le « quitte à » ne couvre
pas, la dualité échangeant $\alpha_0$ et $\alpha_1$ sans changer leurs signes :
$(4, 5')$, $(5', 4)$ et $(5', 5')$, c'est-à-dire $\alpha_1 = c_5' < 0$, ou
$\alpha_0 = \alpha_1 = c_5'$. Ils passent le test de positivité. Ils sont exclus
par le point III), qui place $G_0^+$ parmi les groupes finis de rotations :
aucun ne contient à la fois des éléments d'ordre $4$ et $5$, ce qui écarte
$(4, 5')$ et $(5', 4)$ ; et $\mathfrak{A}_5$ n'a pas de couple générateur formé
de deux éléments d'ordre $5$ de la même classe dont le produit soit d'ordre
$2$, ce qui écarte $(5', 5')$\footnote{Ce complément est de l'édition ; le
dernier fait se lit dans le recensement de la section VI (aucun type
$(5, 5, 2)$ ni $(5', 5', 2)$). La page écrit dans les deux « OPS » un $\pm$ dont
le $+$ est souligné deux fois. La matrice a des coefficients écrits sur
d'autres, biffés ; son déterminant est de l'édition.}. Le théorème est donc
juste.

\pagerange{82}{93}
\section*{IX. Opérations sur les systèmes de Schwarz (pages 82 à 93)}

Cette liasse est la moins lisible du dossier : entre les formules, la prose de
liaison des pages 84 à 89 est en grande partie perdue, et les pages 91 à 93 sont
des brouillons de calcul. On donne la construction, qui se lit, et les
résultats annoncés.

\pagerange{82}{84}
\subsection*{Le groupe $\Pi$ et ses automorphismes (pages 82 à 84)}

Soit $\Pi_0 = \langle \ell_0, \ell_1, \ell_2 \mid \ell_0\ell_1\ell_2 = 1\rangle$,
libre de rang $2$ --- le groupe fondamental de la sphère privée de trois
points --- et $\widehat\Pi$ son complété profini. Un homomorphisme de $\Pi_0$
dans un groupe $H$ (de $\widehat\Pi$ si $H$ est fini) est un triple
$(\rho_0, \rho_1, \rho_2)$ de produit $1$ ; il est surjectif si et seulement si
c'est un système de Schwarz. Soient $\mathcal{T}(H) \subset H^3$ l'ensemble de
ces triples et $\mathrm{Sch}(H)$ celui des systèmes de Schwarz. Les groupes
\[
  \Omega_0 = \mathrm{Aut}(\Pi_0) \longrightarrow \Omega = \mathrm{Aut}(\widehat\Pi)
\]
agissent à droite sur $\mathcal{T}(H)$ et $\mathrm{Sch}(H)$. On a des suites
exactes
\[
  1 \to \Pi_0 \to \Omega_0 \to \widetilde\Omega_0 = \mathrm{Out}(\Pi_0) \to 1,
  \qquad 1 \to \widehat\Pi \to \Omega \to \widetilde\Omega =
  \mathrm{Out}(\widehat\Pi) \to 1 ,
\]
et $\widetilde\Omega_0$, $\widetilde\Omega$ agissent sur $\mathcal{T}(H)/H$ et sur
$\mathrm{Sch}(H)/H$\footnote{La page met un « ? » sur l'injectivité de
$\widehat\Pi \to \Omega$ : elle a lieu, un groupe profini libre de rang $\geq 2$
étant de centre trivial. On sait aussi que $\mathrm{Out}(F_2) \simeq
\mathrm{GL}_2(\mathbf{Z})$ (Nielsen) ; la page n'en dit rien. Les petites
lettres au-dessus du chapeau de $\widehat\Pi$ ne se lisent pas.}.

$\mathfrak{S}_3$ opère fidèlement sur $\Pi_0$ et $\widehat\Pi$ :
$\sigma(\ell_i) = \ell_{\sigma(i)}$ si $\sigma$ est paire, $\ell_{\sigma(i)}^{-1}$
si $\sigma$ est impaire, et l'injection dans $\widetilde\Omega$ « se vérifie
immédiatement » ; c'est l'action de la section VI sur les triples. Soient
$\gamma_0 = \mathbf{Z}^\times = \{\pm 1\}$ et $\gamma = \widehat{\mathbf{Z}}^\times$.
Pour une partie $I$ à deux éléments de $\{0, 1, 2\}$ et $\varepsilon =
(\varepsilon_i)_{i \in I}$, on définit l'automorphisme $\mathsf{e}^I(\varepsilon)$
par $\ell_i \mapsto \ell_i^{\varepsilon_i}$ pour $i \in I$, le troisième générateur
étant déterminé par la relation ; par exemple, pour $I = \{0, 1\}$,
\[
  \ell_0 \mapsto \ell_0^{\varepsilon_0}, \qquad \ell_1 \mapsto
  \ell_1^{\varepsilon_1}, \qquad \ell_2 \mapsto \ell_1^{-\varepsilon_1}
  \ell_0^{-\varepsilon_0} .
\]
Cela définit des homomorphismes $\mathsf{e}^I : \gamma_0^I \to \Omega_0$,
$\gamma^I \to \Omega$\footnote{Il les note $\rho_0^I$, $\rho^I$, et, pour les trois
$I$, $\rho_{\varepsilon,\varepsilon'\cdot}$, $\rho_{\cdot\varepsilon\varepsilon'}$,
$\rho_{\varepsilon\cdot\varepsilon'}$, le point marquant le générateur sur
lequel l'opération n'agit pas par une puissance ; la place du point se lit
mal. On écrit $\mathsf{e}$ pour garder $\rho$ aux rotations. À la deuxième ligne de
la page 83 on lit $\ell_2^{-\varepsilon'}\ell_2^{-\varepsilon}$, où il faut
$\ell_2^{-\varepsilon'}\ell_1^{-\varepsilon}$. Pour $\varepsilon \in
\widehat{\mathbf{Z}}^\times$ ce sont bien des automorphismes, un endomorphisme
surjectif d'un groupe profini de type fini étant bijectif.}, et les relations
\[
  \sigma\, \mathsf{e}^I(\varepsilon)\, \sigma^{-1} = \mathsf{e}^{\sigma I}(\sigma
  \varepsilon) \qquad (\sigma \in \mathfrak{S}_3) .
\]
Le sous-groupe $\Omega^+$ (resp. $\Omega_0^+$, et leurs images extérieures)
engendré par $\mathfrak{S}_3$ et les $\mathsf{e}^I(\gamma^I)$ contient le
sous-groupe distingué $\mathbb{E}\mathrm{l}$ engendré par les seuls
$\mathsf{e}^I$ --- le « groupe des éliminations » --- et
\[
  \Omega^+ \simeq \mathfrak{S}_3 \ltimes \mathbb{E}\mathrm{l}, \qquad \Omega_0^+
  \simeq \mathfrak{S}_3 \ltimes \mathbb{E}\mathrm{l}_0 ,
\]
et de même pour les groupes extérieurs\footnote{Que le produit soit
semi-direct, c'est-à-dire $\mathfrak{S}_3 \cap \mathbb{E}\mathrm{l} = 1$, est
affirmé dans une phrase dont presque tous les mots manquent ; on rapporte la
conclusion. « Éliminations » est la lecture de la transcription, le mot est
entre guillemets sur la page.}.

\pagerange{84}{88}
\subsection*{Une présentation (pages 84 à 88)}

\textbf{Construction} (pages 84 et 85). Pour deux groupes $\mathfrak{S}$ et
$\gamma$, le noyau $\mathcal{E}$ de $\mathfrak{S} * \gamma \to \mathfrak{S}$ est
le produit libre des conjugués $\sigma\gamma\sigma^{-1}$, $\sigma \in
\mathfrak{S}$ :
\[
  \mathfrak{S} * \gamma \simeq \mathfrak{S} \ltimes \gamma^{(*\mathfrak{S})} ,
\]
ce qu'on voit par la propriété universelle. Avec $\mathfrak{S} = \mathfrak{S}_3$
et $\gamma = \widehat{\mathbf{Z}}^\times$, on obtient un homomorphisme
$\mathfrak{S}_3 * \gamma \to \Omega$ ; il n'est pas injectif, car
$\mathsf{e}_{\varepsilon,1}$ et $\mathsf{e}_{1,\varepsilon'} = \pi\,
\mathsf{e}_{\varepsilon,1}\pi^{-1}$ ($\pi = \pi_{01}$) commutent. En imposant ces
commutations on obtient
\[
  \mathfrak{S}_3 \ltimes (\gamma \times \gamma)^{(*\, \mathfrak{S}_3/\pi_{01})}
  \longrightarrow \Omega^+ \subset \Omega ,
\]
produit libre de trois copies de $\gamma \times \gamma$ permutées
transitivement par $\mathfrak{S}_3$, et de même avec $\gamma_0$ pour
$\Omega_0^+$. « Il serait intéressant d'examiner les propriétés d'injectivité et
surjectivité » (page 86) ; l'image de $\mathcal{E}$ est $\mathbb{E}\mathrm{l}$
(resp. $\mathbb{E}\mathrm{l}_0$), et la question « $\Omega^+ = \Omega$ ? » reste
posée (page 87)\footnote{La ligne de la page 86 qui énonce les relations est
serrée et peu lisible. « $\gamma_0 = \{\pm1\}$ » est noté dans un coin.}.

\textbf{Réduction à l'exposant $6$} (pages 87 et 88). Soit $\Pi(6)$ le quotient de
$\Pi$ par les relations $g^6 = 1$, $\Omega(6) = \mathrm{Aut}(\Pi(6))$. La
surjection $\widehat{\mathbf{Z}}^\times \to (\mathbf{Z}/6)^\times \simeq \{\pm 1\}$
a une section $\gamma_0 \hookrightarrow \gamma$, et le composé
$\mathfrak{S}_3 \ltimes (\gamma \times \gamma)^{(*3)} \to \Omega \to \Omega(6)$
se factorise par $\mathfrak{S}_3 \ltimes (\gamma_0 \times \gamma_0)^{(*3)}$. Donc,
pour un groupe $H$ d'exposant divisant $6$ --- par exemple $\mathfrak{A}_4$ ---,
l'action de $\Omega^+$ sur $\mathcal{T}(H)$ et $\mathrm{Sch}(H)$ se fait à
travers ce groupe discret\footnote{Pages 87 et 88, la prose est presque
entièrement illisible ; on rapporte ce que disent les diagrammes et la phrase
finale de la page 88. La transcription écrit d'abord $\overline\Pi(6)$, puis
$\Pi(6)$.}.

\pagerange{89}{90}
\subsection*{Les $96$ systèmes du tétraèdre ; l'exposant $N$ (pages 89 et 90)}

Le groupe $\mathfrak{S}_3 \ltimes (\gamma_0 \times \gamma_0)^{(*3)}$ est engendré
par $\mathfrak{S}_3$ et six involutions $\mathsf{e}_\sigma$ ($\sigma \in
\mathfrak{S}_3$), avec les relations
\[
  \mathsf{e}_\sigma = \sigma\, \mathsf{e}_1 \sigma^{-1}, \qquad \mathsf{e}_1^2 = 1,
  \qquad \mathsf{e}_1 \mathsf{e}_\pi = \mathsf{e}_\pi \mathsf{e}_1 \quad (\pi =
  \pi_{01}) ,
\]
d'où $\mathsf{e}_{\tau\sigma} = \tau\mathsf{e}_\sigma\tau^{-1}$, $\mathsf{e}_\sigma^2
= 1$ et $\mathsf{e}_\sigma\mathsf{e}_{\sigma\pi} = \mathsf{e}_{\sigma\pi}
\mathsf{e}_\sigma$. Ce groupe des « opérations élémentaires » agit sur les
$72 + 24 = 96$ systèmes de Schwarz de $H \simeq \mathfrak{A}_4$ (section VI), et
sur $\mathrm{Sch}(H)/H$, qui a $8$ éléments : $(3,3,2)^+$, $(3,3,2)^-$ et leurs
transformés par $\mathfrak{S}_3^+$, $(3,3,3)^+$ et $(3,3,3)^-$ ; modulo
$\widetilde H = \mathrm{Aut}(H) \simeq \mathfrak{S}_4$, il en reste $4$, de types
$(3,3,2)$, $(3,2,3)$, $(2,3,3)$, $(3,3,3)$. « Il serait intéressant d'expliciter
les groupes quotients finis de $\mathfrak{S}_3 \cdot \mathcal{E}_0$ qui » agissent
sur ces ensembles\footnote{Les nombres $96$, $8$ et $4$ sont confirmés par
énumération ; l'exposant $\pm$ distingue systèmes canoniques et
anticanoniques. La page s'arrête sur des points de suspension.}.

\emph{NB} (page 90). Pour les groupes $H$ dont tous les éléments ont un ordre
divisant $N$ --- $N = 6$ pour le groupe du tétraèdre, $12$ pour
$\mathfrak{S}_4$, $30$ pour $\mathfrak{A}_5$ --- on remplace $\Pi$ par $\Pi(N)$,
$\Omega$ par $\Omega(N)$, et $\mathfrak{S}_3 \ltimes \mathbb{E}\mathrm{l}$ par un
quotient de $\mathfrak{S}_3 \ltimes (\gamma(N) \times \gamma(N))^{(*3)}$, où
$\gamma(N) = (\mathbf{Z}/N)^\times$, qui vaut $\{\pm1\}$, $\{\pm1\}^2$,
$\{\pm1\} \times \mathbf{Z}/4$, et $\{\pm 1\}^2 \times \mathbf{Z}/4$ pour
$\mathfrak{S}_5$ ($N = 60$).

On reconnaît aujourd'hui dans ces opérations l'action de $\mathrm{Aut}(F_2)$
sur les couples générateurs d'un groupe fini --- les classes de Nielsen, ou
« $T$-systèmes » de B. H. et H. Neumann (1951) ---, et, dans les opérations
$\ell_i \mapsto \ell_i^{\varepsilon}$, $\varepsilon \in \widehat{\mathbf{Z}}^\times$,
la forme qu'a l'action de $\mathrm{Gal}(\overline{\mathbf{Q}}/\mathbf{Q})$ sur le
groupe fondamental profini de $\mathbb{P}^1 \setminus \{0, 1, \infty\}$, à
conjugaison près des générateurs, par le caractère cyclotomique : c'est la
rubrique (4) du plan de la page 1\footnote{Le rapprochement est de l'édition.
L'action galoisienne envoie chaque $\ell_i$ sur un \emph{conjugué} de
$\ell_i^{\chi(g)}$, non sur $\ell_i^{\chi(g)}$ lui-même, et ne se réduit donc pas
aux $\mathsf{e}^I$ ; le groupe de Grothendieck--Teichmüller (Drinfeld, 1990)
est postérieur. Ces pages ne disent pas ce rapport.}.

\pagerange{91}{93}
\subsection*{Brouillons (pages 91 à 93)}

Les trois dernières pages de la liasse sont des calculs sans phrases. Page 91,
des opérations élémentaires $\eta_1, \eta_2, \eta'_1, \ldots, \eta''_2$ sur
$(\ell_0, \ell_1, \ell_2)$ --- par exemple $\eta_1 : (\ell_0, \ell_1, \ell_2) \mapsto
(\ell_0^{-1}, \ell_1, \ell_1^{-1}\ell_0)$ --- et le composé $\mathsf{e}_1 =
\eta_1\eta''_2$, qui envoie $\ell_0 \mapsto \ell_0$, $\ell_1 \mapsto
\ell_0^{-2}\ell_1$, $\ell_2 \mapsto \ell_2\ell_0^2$ ; puis un brouillon d'algèbre
linéaire sur $\mathbf{Z}^3$ (l'abélianisé). Pages 92 et 93, les systèmes du
tétraèdre étiquetés par les repères --- $\Sigma_r = (\rho_f, \rho_s, \rho_a)$ de
type $(3,3,2)^+$ pour $r$ direct, $(3,3,2)^-$ pour $r$ inverse, $(\rho_s,
\rho_{s'}, \rho_{s''})$ de type $(3,3,3)^\pm$ --- et l'effet des opérations
$\mathsf{e}_{-1,1}$, $\mathsf{e}_{1,-1}$ : elles échangent $(3,3,2)^+$ et
$(3,3,3)^+$, $(3,3,2)^-$ et $(3,3,3)^-$, ou croisent les signes ; enfin
$\mathfrak{S}_3 * \mathbf{Z}/2 \simeq \mathfrak{S}_3 \ltimes
(\mathbf{Z}/2)^{(*\mathfrak{S}_3)}$, et la remarque que $\mathsf{e}(\varepsilon)$
commute à son conjugué par une transposition\footnote{Il note les opérations de
la page 91 d'un $\xi$ bouclé, distinct de son $\rho$ ; on écrit $\eta$, $\xi$ étant
réservé à la section VII. Les signes de plusieurs triples de la page 92 sont
surchargés, les exposants de $\ell_2$ de la page 91 aussi ; les flèches du
tableau sont dessinées de droite à gauche. La page 92 porte un tétraèdre
($s$ en haut, $s'$ à droite, une face hachurée), la page 93 un autre ($s$, $s'$,
$a$, $f$). On ne donne de ces brouillons que ce qui se lit sûrement.}.

\pagerange{95}{97}
\section*{X. Classes de sous-groupes de $\mathfrak{S}_5$ (pages 95 à 97)}

Soit $T$ un ensemble à cinq éléments, $\mathfrak{S}_5 = \mathfrak{S}_T$. La
page 95 range les sous-groupes selon la figure de $T$ qu'ils stabilisent : aux
triangles de $T$ (parties à trois éléments, $10$) correspondent les classes
$\mathbf{Z}_3 = \mathfrak{A}_3$, $\mathbf{Z}_6$, $\mathfrak{S}_2$, $\mathfrak{S}_3$,
$\mathbb{D}_3$ (le $\mathfrak{S}_3$ « tordu » $\langle (abc), (ab)(de)\rangle$),
$\mathbb{D}_6 = \mathfrak{S}_3 \times \mathfrak{S}_2$ ; aux « carrés » ($15$),
$\mathbf{Z}_2$ (doubles transpositions), $\mathbf{Z}_4$, $\mathfrak{S}_{2,2}$,
$\mathbb{D}_4$ ; aux points ($5$), $\mathbb{D}_2$, $\mathfrak{A}_4$,
$\mathfrak{S}_4$ ; aux pentagones, $\mathbf{Z}_5$ et $\mathbb{D}_5$ ; et $1$,
$\mathfrak{A}_5$, $\mathfrak{S}_5$ sont distingués. D'où le compte
\[
  6 \times 10 + 4 \times 15 + 3 \times 5 + 2 \times 12 + 3 \times 1 = 162 .
\]
Page 97 : « les 18 classes de sous-groupes de $\mathfrak{S}_5$, dont 9
abéliens » ; $9$ classes de sous-groupes de $\mathfrak{A}_5$, $11$ de
$\mathfrak{S}_4$, $5$ de $\mathfrak{A}_4$, $4$ de $\mathfrak{S}_3$, $7$ de
$\mathbb{D}_4$, $10$ de $\widetilde{\mathfrak{S}}_3 = \mathfrak{S}_3 \times
\mathfrak{S}_2$ ; et des identités : $\mathbf{F}_2 = \mathfrak{Z}(\mathbb{D}_4)$,
$\mathbb{D}_4 = \mathcal{N}_{\mathfrak{S}_4}(\mathbf{F}_2) =
\mathcal{N}_{\mathfrak{S}_4}(\mathbf{Z}_4)$, $\mathbf{F}_2 = {}_2\mathbf{Z}_4$,
$\mathfrak{S}_4 = \mathcal{N}_{\mathfrak{S}_5}(\mathbf{F}_2^2)$, $\mathbf{F}_2^2$
étant le seul sous-groupe distingué d'indice $6$ de $\mathfrak{S}_4$.

Les identités sont justes, et les nombres $9$, $11$, $5$, $4$, $10$ aussi. Mais
$\mathfrak{S}_5$ a $156$ sous-groupes en $19$ classes, non $162$ en $18$ : il
manque le groupe de Frobenius $F_{20} = \mathrm{AGL}(1, \mathbf{F}_5)$,
normalisateur d'un $5$-Sylow ($6$ conjugués), et les pentagones ne
correspondent pas bijectivement aux sous-groupes : chaque $\mathbf{Z}_5$ (et
chaque $\mathbb{D}_5$) stabilise deux des douze pentagones, le pentagone et le
pentagramme qu'il dessine. Le compte juste est $60 + 60 + 15 + 3 \times 6 + 3
= 156$ ; les « 9 abéliens » sont exacts (trivial compris), et $\mathbb{D}_4$ a
$8$ classes de sous-groupes, non $7$\footnote{Nombres recalculés par
énumération pour cette édition. « $\mathfrak{S}_{22}$ » est le
$\langle (ab), (cd)\rangle$ non distingué. La page 95 note les groupes alternés
d'une lettre en forme de $G$. La page 96 ne porte que des dessins au crayon
(une bipyramide triangulaire, un assemblage de tétraèdres autour d'un
sommet) ; les treillis des pages 95 et 97, qui mettent au net les inclusions
entre classes, ne sont pas redessinés. La note marginale de la page 95 se
propose de « déterminer les automorphismes extérieurs des groupes
$\mathfrak{S}_n$ ».}.

\end{document}
